% =====================================================================
%  R. Nielsen, OM SKJÆBNE OG FORSYN.  Kjøbenhavn: Otto Schwarz, 1853
%  (Berlingske Bogtrykkeri).  83 pp.  Fraktur.  Reissued unchanged 1867 as
%  »Om Tilværelsens ideale Magter: Skjæbne og Forsyn« (83 pp.).
%
%  TRANSCRIPTION -- GENERATED; DO NOT EDIT BY HAND.  Rebuild, chained (&&), from the repo root:
%    python3 ebook-method/skjaebne-og-forsyn/draft.py --emit --force
%      this template + the scan's text layer (Google's OCR in the Google Books scan of the
%      British Library copy, books.google.com/books?id=tHRZAAAAcAAJ;
%      ~/bibliotek/Nielsen, Rasmus/1853-skjaebne-og-forsyn.pdf, see SCANS.tsv);
%    python3 ebook-method/collate/check.py skjaebne-og-forsyn --witness ocr --apply
%      the decisions in decisions.tsv, taken against tesseract's independent reading of the images;
%    python3 ebook-method/skjaebne-og-forsyn/oepass.py --apply
%      the ø pass (lexicon + 147 tight word crops read at the image, oe-reads.tsv);
%    python3 ebook-method/skjaebne-og-forsyn/finish.py --apply
%      headings, the footnote, divided words, emphasis, PRINTED AS IS marks (each from a reading at the image).
%
%  WORD-ACCURACY (2026-10-07; IN PROGRESS).  No e-book, no earlier transcription exists.  The draft (Google's OCR
%  layer) was collated word by word, with punctuation and word division, against tesseract (tessdata_best
%  script/Fraktur), with deu_latf and frak2021 as third and fourth readings: 1264 disagreements.  198 settled by
%  check.py's own rules, 786 by rules.py (each decision named; a blind audit of 46 of them, mixed into the crops,
%  found no rule error); the rest at the image, in BLIND crops read by three Sonnet readers who saw neither OCR
%  reading (6 of 6 planted controls read right), and a few by hand from the readers' notes.  The ø (read o by every
%  OCR): lexicon + 147 tight word crops (for/før ...).  Title page, headings, the one footnote (p. 50), the page
%  heads of pp. 3 and 5 and the residue were read at the image (structure reader); emphasis (letter-spacing, bold)
%  from 80 word crops.
%
%  SWEEP (2026-10-07): every word of the text not attested elsewhere in the corpus (97, proper names aside) read at
%  the image in a crop: 8 mended (word-fixes.tsv: Ære for the layer's „Wre“ x3, Lysenes, gjælder, timelig,
%  personlige, og), the rest confirmed as printed.  Emphasis: every span re-read as a whole line (r-sweep.tsv L lines).
%
%  SAMPLED PAGE CHECK (2026-10-07, pagecheck.py; seed 20261007: pp. 10, 29, 42, 50, 57, 60, 67, 1177 words, transcribed
%  BLIND at the image and compared word by word, every difference settled): 5 word errors (4.25 / 1000, 95 % upper 8.9)
%  and 2 of punctuation.  Four of the five belonged to two classes of fault that the layer AND tesseract share, so the
%  collation could not see them: the Fraktur Ø read as D (Øie, Øieblik, Øiekast, Øine) and j read as i after g/k (gjør,
%  afgjørende, kjende, skjøndt, gjengjældende).  Every occurrence of both classes in the book (49) was looked at in a line
%  crop and mended (finish.py CLASS; Forskiellen p. 8 and Kiød p. 82 are printed with i).  The fifth: sheet signatures
%  (1*, 2*, 4*, 5*) had been kept in the text; all dropped.  Punctuation: „Lidenskaben?“ (p. 10) and four more marks lost
%  where a blind reading's box stopped short of them, restored; „men — dersom“ (p. 42).
%
%  SECOND READING BY TRANSLATION (2026-10-07).  Two translators listed 36 odd spots; all were read at the image in line
%  crops (main session).  15 transcription faults mended (finish.py TRFIX): a page's lost first word („Middel-|maadighed“,
%  p. 11), specks read as marks (Det,. p. 7; De -moralske p. 10; skal. p. 82; In-. p. 80), „j“ for i (p. 34), straight
%  quotes (p. 36), a stray closing quote (p. 54), a false paragraph break inside a quotation (p. 64), letere/lettere
%  (p. 74), gage/gaae (p. 82); and two words the rule `tokens` had glued on letter-spaced lines (EdersFiender p. 29,
%  personligeOpgaver p. 58), found by prescreen.py before dispatch.  Six more misprints confirmed and marked (Jorgens,
%  Personligligheden, Ham.,, hiælper, Nøst, deri,.).
%
%  NOT ESTABLISHED (yet).  EMPHASIS is the least certain layer: the print letter-spaces its Scripture quotations
%  throughout, and the extent of the runs on pp. 29 (Eder: „elsker Eders Fiender …“), 42 („Ja endog alle Eders
%  Hovedhaar …“ -- since settled), 55-58 („den, som giør fattig …“) has not been set; on p. 48 („det umuligt, men for Gud
%  ere alle Ting mulige“) the line reader saw bold where the word reader saw spacing; „hvorledes“ (p. 4) and „ikke“
%  (p. 30) are set on one reading against another.  (The p. 42 quotation is now spaced throughout, seen at the image.)  Open: „op voxte“ across the page turn pp. 12-13 (one word?);
%  „paà“ (p. 20, a grave-like mark, kept as the layer read it).  No whole page has yet been read blind against the text.
%  PRINTED AS IS: „ørgeligt“ for sørgeligt (p. 18); „den Rige. en“ (p. 56).
%
%  PAGE MAP.  printed p = PDF page - 2 throughout (p. 3 = PDF 5, p. 83 = PDF 85); folios of 79 pp. read
%  off the layer, all consistent (to be confirmed at the image).  PDF 3: title, p. [1]; PDF 4: p. [2];
%  PDF 86: the British Library's stamp; PDF 88: Otto Schwartz's advertisement (not transcribed).
%
% =====================================================================
\documentclass[12pt,a4paper,openany]{book}
\usepackage[utf8]{inputenc}
\usepackage[T1]{fontenc}
\usepackage{libertinus}
\usepackage{graphicx}
\DeclareUnicodeCharacter{0254}{\reflectbox{c}}  % ɔ (U+0254) = reversed c
\DeclareUnicodeCharacter{A75B}{r}   % r rotunda (ꝛc. = etc.)
\usepackage[danish]{babel}
\usepackage[protrusion=true,expansion=true]{microtype}
\usepackage{setspace}
\usepackage[margin=1.3in]{geometry}
\usepackage{fancyhdr}
\usepackage{hyperref}
\hypersetup{pdftitle={Om Skjæbne og Forsyn}, pdfauthor={Rasmus Nielsen}, hidelinks}
\setlength{\headheight}{15pt}
\pagestyle{fancy}
\fancyhf{}
\fancyhead[LE,RO]{\thepage}
\fancyhead[C]{\textit{R. Nielsen: Om Skjæbne og Forsyn (1853)}}
\renewcommand{\thefootnote}{\ifcase\value{footnote}\or *)\or **)\or ***)\else
  \arabic{footnote})\fi}
\setstretch{1.3}
\usepackage{marginnote}
\newcommand{\opage}[1]{\marginnote{\footnotesize\textit{[#1]}}}
\newcommand{\apage}[1]{\marginnote{\footnotesize\textit{[#1]}}}
% headings as printed (read at the image, .parts/notes/r-struct.tsv): the Indledning and the part titles in large
% Fraktur, the part letters A./B./C. and section numerals I.-IV. in bold antiqua, section titles in bold Fraktur; a short
% rule under each part title; in B, sections I and II take a bold Scripture verse as their title (\motto)
\newcommand{\overskrift}[1]{\begin{center}{\Large #1}\\[0.3em]\rule{3em}{0.4pt}\end{center}\addcontentsline{toc}{section}{#1}}
\newcommand{\afdeling}[3][]{\begin{center}\ifx&#1&\else#1\\[0.6em]\fi\textbf{#2}\\[0.3em]{\Large #3}\\[0.3em]\rule{3em}{0.4pt}\end{center}\addcontentsline{toc}{section}{#2 #3}}
\newcommand{\delafsnit}[1]{\begin{center}{\Large\textbf{#1}}\end{center}\addcontentsline{toc}{subsection}{#1}}
\newcommand{\stykke}[2]{\begin{center}\ifx&#1&\textbf{#2}\else\textbf{#1}\ifx&#2&\else\\[0.2em]\textbf{#2}\fi\fi\end{center}}
\newcommand{\motto}[1]{\begin{quote}\textbf{#1}\end{quote}}

\begin{document}

% --- PDF 3, the title page, p. [1], read at the image 2026-10-07 (.parts/notes/r-struct.tsv, F3): Fraktur
% throughout but the year (antiqua); an ornamental rule (drawn plain) between the author and the imprint.  The print has
% „Schwartz“ (the layer read Schwarz; KB's catalogue has Schwarz).  The verso, p. [2] (PDF 4), is blank.
\begin{titlepage}
\begin{center}
\vspace*{4em}
{\Large Om}\\[1em]
{\Huge Skjæbne og Forsyn}\\[1.5em]
af\\[1em]
{\Large\textbf{R. Nielsen,}}\\[0.5em]
{\small Professor i Philosophi ved Kjøbenhavns Universitet.}\\[3em]
\rule{8em}{0.4pt}\\[3em]
\textbf{Kjøbenhavn.}\\[0.5em]
Forlagt af Otto Schwartz.\\[0.5em]
{\small Berlingske Bogtrykkeri.}\\[0.5em]
\textit{1853.}
\end{center}
\end{titlepage}

% ===== TEXT (pp. 3--83) =====
% --- p. 3 ---
\setcounter{footnote}{0}%
\opage{3}\overskrift{Indledning.}

Det er umuligt ved en almindelig Betragtning af de
historiske Begivenheder at erkjende Forskjellen mellem
Skjæbne og Forsyn.

Da Jerusalem blev forstyrret, da sagde Jøden:
„det er Jehova, der hævner sin Lovs Overtrædelse“;
da sagde Nazaræeren: „det er den Korsfæstedes Blod,
der kommer over Folkets Hoved“; men Romeren sagde:
„det er Skjæbnen.“

Som nu Jøden, Nazaræeren og Romeren dengang
dømte, hver efter sin Stemning og sit Standpunkt,
saa dømmer endnu den Dag idag hvert Parti og
hvert Menneske efter sin Stemning og sit Standpunkt
om enhver historisk Begivenhed. Det kommer altsaa
ikke an paa Begivenheden, men paa Stemningen og
Standpunktet.

Det er umuligt ved en almindelig Betragtning af
de enkelte mærkelige Tildragelser i de enkelte Menneskers Liv at gjøre Forskjel imellem Skjæbne og Forsyn.

To Mennesker kunne i det Udvortes opleve de
samme Hændelser, komme i de samme Farer,  frelses
% --- p. 4 ---
\setcounter{footnote}{0}%
\opage{4}paa samme uventede og vidunderlige Maade: den Ene
seer i alt Dette et Skjæbnens Spil, den Anden seer
et Forsyn. Det kommer altsaa slet ikke an paa,
hvad det er, der opleves, men \emph{hvorledes} det opleves;
det kommer ikke an paa Tildragelsen, men paa
Tænkemaaden.

Det er umuligt for et Menneske at faae sin Tro
paa Forsynet ligefrem stadfæstet ved de Begivenheder,
der møde ham; thi han maa jo tyde alle Begivenheder
efter sin Tro. Kommer Medgang: det er Forsynet;
kommer Modgang: det er ogsaa Forsynet; er Alt i
Tilværelsen forstaaeligt og lykkeligt, og Fremtiden lys:
det er Forsynet; er Alt i Tilværelsen forvirret og
besværligt, og Fremtiden mørk: det er ogsaa Forsynet.

Den, der ikke saaledes i Alt underordner sin egen
menneskelige Dom under sin Tro paa Forsynet, troer
slet ikke paa Forsynet.

Men, naar Begivenhederne saavel i Historien som
i de enkelte Menneskers Liv lade sig udtyde paa saa
høist forskjellig Maade: saa kunne vi jo, hvad Skjæbne
og Forsyn angaaer, Intet, aldeles Intet lære af Begivenhedernes
Gang.

% --- p. 5 ---
\setcounter{footnote}{0}%
\opage{5}\afdeling{A.}{Om Skjæbnen.}
\stykke{I.}{Hvad er Skjæbnen?}

Den kolde, ubøielige og dog saa lunefulde,
vilkaarlige Magt, vi kalde Skjæbne, er væsentlig
begrundet i Naturtingenes indbyrdes Sammenhæng,
i Tilværelsens lovbestemte Orden,
altsaa i Verdensfornuften selv.

Fornuftloven og den ved samme betingede
Nødvendighed er imidlertid i og for sig endnu
ingenlunde at kalde Skjæbne. Forsaavidt de
tilværende Ting umiddelbart underordne sig de
almeengyldige Love, er Nødvendigheden, den
ydre saavelsom den indre, at betragte som Vexelvirkningens,
Heelhedens og Harmoniens Grundlag.
Stjernernes Gang, Aarstidernes Vexel,
% --- p. 6 ---
\setcounter{footnote}{0}%
\opage{6}Planternes Væxt, alt Sligt er begrundet i Natursammenhængen,
og forsaavidt i Nødvendigheden;
men her er ingen Skjæbne.

Først i Dyreverdenen, hvor det individuelle
Liv yttrer sig med Selvfornemmelse og selvisk
Formaal, begynder Skjæbnen at dæmre. Hvad
kan et Dyr ikke lide; hvor forskjellig er ikke i
denne Henseende Dyrenes Skjæbne: den fangne
Fugl og den frie Fugl have hver sin Skjæbne;
og nu Huusdyrenes Skjæbne, --- dersom Heste
og Æsler kunde tale!

Men Heste og Æsler kunne ikke tale; denne
Ufuldkommenhed mildner deres Skjæbne. Dyrets
Afhængighed af Instinctet, dets Fortabthed
i Øieblikket, dets Bevidstheds Dunkelhed, der
tilslører Fremtidens Muligheder og skjuler det
Forbigangne i Glemsel, forhindrer det fra at
sætte sig som fri Selvhensigt, bøier det mildt
under Naturnødvendighedens Aag, og er som
et Dække, der afbøder Skjæbnens Slag.

Kun hvor der er en Selvhensigt, en Villiens
Plan, et personligt Maal, kan der være
Tale om en Magt, der kuer Hensigten, hindrer
Planen og forstyrrer Maalet. Blandt alle Jorgens Skabninger er derfor Mennesket den eneste,  % PRINTED AS IS: p. 6 „Jorgens“ for Jordens (seen at the image)
% --- p. 7 ---
\setcounter{footnote}{0}%
\opage{7}hvem Naturen har indrømmet det store, tvetydige
Fortrin, at han i egentlig Forstand, saalænge
han staaer paa Verdslighedens Grund,
skal staae under Skjæbnen. Rige Anlæg, mægtig
Villie, dristig Aand, alle Personlighedens
naturlige Gaver kunne ikke frelse ham fra den
mørke, gaadefulde Magt; thi jo større Personligheden
er, desto mægtigere er Skjæbnen.
Skjæbnen er alt det Personliges Fiende; derfor
maae Guderne føle deres Magt, derfor maa
Jupiter med sin Lynstraale bøie sig under Fatum.

\stykke{II.}{Skjæbnens Herredømme.}

Er det et Skjæbnens Særkjende, at den
overalt træder Personligheden fiendtlig imøde;
saa følger jo heraf, at dens Herredømme
maa være ligesaa udstrakt, som den personlige
Stræben.

I de timelige Goders Erhvervelse, i Ideernes
Sphære, ja selv i det Moralske, maae vi
skjelne imellem Det, der beroer paa os, og
Det, der ikke beroer paa os. Dersom nu Det,
der ikke beroer paa os, tog Hensyn til vor
personlige Trang, og kom vor fornuftige Villie
% --- p. 8 ---
\setcounter{footnote}{0}%
\opage{8}til Hjælp, saa var det ikke Skjæbnen, men en
høiere, Alt styrende Viisdom. Verdslig seet,
er det Omvendte just Tilfældet. Den almene
Overmagt, hvoraf vor Stræben afhænger, er
saa langt fra at ændse det Personlige i os, at
den meget mere forholder sig ligegyldig, ja
endog spottende dertil: dette er just det Fatale.

Betragte vi Forskiellen mellem Sunde og
Syge, Rige og Fattige, Høie og Ringe i Verden,
da indsees det ved første Øiekast, at Fordelingen
af de timelige Goder ingenlunde skeer
efter Trang og Fortjeneste, men at den fordelende Verdensmagt har fulgt Vilkaarligheden,
har ladet Held og Uheld, Lykke og Ulykke raade
for Delingen, og saavel i Skikkelse af Lykke som
Ulykke havt den menneskelige Personlighed til
Bedste.

Vende vi os fra det Udvortes til det Indvortes, da gjentager sig en lignende Vilkaarlighed.
Det er et menneskeligt Fortrin at kunne
hæve sig over det Partikulære, det Smaalige,
sætte sig almeengyldige Formaal, det vil sige:
leve for Ideer. Men dette Fortrin er ikke
ligeligt for Alle; her møde vi jo den skjæbnesvangre
Forskjel imellem de Begavede og de
% --- p. 9 ---
\setcounter{footnote}{0}%
\opage{9}Ubegavede, de Aandrige og de Aandløse. Det
Udtryk, at leve for en Idee, er desuden tvetydigt,
saalænge det lades ubestemt, til hvilken
Sphære Ideen hører. Der gives personlige
og upersonlige Ideer.

Kunstens, Videnskabens, Statens Ideer ere
upersonlige; Moralens Idee er personlig.

At leve for en upersonlig Idee er at give
sig hen i det Upersonlige. Enhver aandig Hengivelse
i det Aandige forhøier Selvfølelsen; derfor
have Genierne og de Begeistrede en saa kraftig
intensiv Leven i sig selv. Men det Upersonlige
faaer tilsidst Overmagten. Det Inderlighedens
Baand, der binder de Udmærkede til Ideen og
gjør dem til dens kraftige Redskaber, er tillige
den Lænke, hvori Skjæbnen fanger Personligligheden  % PRINTED AS IS: p. 9 „Personlig-|ligheden“, a syllable set twice at the line end
og binder den til sit Hjul: Værket
fuldendes, Mesteren gaaer tilgrunde; Ideen
seirer, dens Redskab ligger sønderbrudt (en
Columbus, en Kepler, en Mozart, en Fulton
kunne, hver i sit Sprog, afgive Forklaring).
Haard er Skjæbnen, naar den vilkaarlig fordeler
de endelige Goder, leger med de Lykkelige
og plager de Ulykkelige; men grusom er den,
naar den som upersonlig Idee lister sig ind
% --- p. 10 ---
\setcounter{footnote}{0}%
\opage{10}i Menneskets Indre og blodsuger Personligheden.

For at modstaae Skjæbnen bør altsaa Mennesket
give sig ind under den moralske Verdensorden
og øve sig i Dyd, thi Dydens Idee er
personlig. Dette Raad er gammelt og godt,
men, verdslig forstaaet, utilstrækkeligt; saa langt,
som Naturlovene gaae, gaaer ogsaa Skjæbnens
Rige.

Som vi med Hensyn paa de upersonlige
Ideer skjelne imellem Begavede og Ubegavede,
saaledes maae vi med Hensyn paa den moralske
Idee gjøre Forskjel imellem ædle og uædle Naturer.
De moralske Naturanlæg synes nemlig
at være ligesaa vilkaarlig fordeelte, som Talenterne
og de timelige Goder. Der gives Individer,
der ikke kunne Andet end være fromme,
dydige, elskelige, fordi det ligger i deres Natur;
og der gives Individer, der af Naturen
ere saa godt som prædestinerede til Forbrydelse.
Vistnok har Mennesket, hvad man kalder en fri
Villie; men hvad formaaer Villien imod Lidenskaben?
hvad bliver det i Virkeligheden til
med den moralske Villie, hvor Naturgrunden
mangler? Halvhed, Hykleri, Qvakleri, Middel%
% --- p. 11 ---
\setcounter{footnote}{0}%
\opage{11}maadighed! Villien griber først kraftig, naar den
selv kraftig er greben, men den gribes kun kraftig
af Det, som opgaaer for den indre Sands,
og tilfredsstiller personlig. Har nu den moralske
Idee alene sin Gyldighed i den personlige
Selvtilfredsstillelse, saa er jo Den, der
tilfredsstiller sig i det Slette, naar han virkelig
tilfredsstiller sig, ligesaa berettiget, som Den,
der tilfredsstiller sig i det Gode; Forskjellen
sættes af Skjæbnen og udslettes igjen af Skjæbnen;
thi Døden gjør os Alle lige: Døden er
den sidste Skjæbne.

\stykke{III.}{Kamp mod Skjæbnen.}

Den verdslige Almeenaand er en upersonlig
Aand, der vil fortære det menneskelige Selv;
den verdslige Personlighed er i sin Grund en
egoistisk Personlighed, som til Trods for det
Almene vil hævde det menneskelige Selv. Imellem
det upersonlige Almene og det egoistisk
Personlige er Forholdet spændt. Det Upersonliges
Overmagt i Forholdet er Skjæbnen; derfor
strider Personligheden imod, og holder sig,
saavidt den formaaer, i Kamp mod Skjæbnen.

% --- p. 12 ---
\setcounter{footnote}{0}%
\opage{12}Efter den særegne personlige Fordring og
den Maade, hvorpaa Fordringen drives igjennem,
charakterisere vi Kampen som: den \emph{borgerlige, den heroiske og den moralske}.

Den borgerlige Verden er Endelighedens,
de timelige Goders Verden. Den fornuftige
Villie kan da i en Verden, der er bygget paa
lutter Betingelser, ikke tænke paa at gjøre
nogen ubetinget Fordring eller kræve ubetinget
Tilfredsstillelse. Det endelige Maal er en taalelig
Tilstand: Personligheden bukker for Omstændighederne, og gaaer paa Accord med Skjæbnen.
Men den Overeenskomst, der saaledes sluttes
mellem det Egoistiske og det Almene, er alligevel
forbigaaende, upaalidelig og falsk; de contraherende
Parter ere væsentlig hinandens Fiender:
den taalelige Tilstand bliver snart utaalelig, og
Personligheden drives til Fortvivlelse. Den
Mægtige trykker den Ringe, og den Ringe hader
den Mægtige; den Rige foragter den Fattige,
og den Fattige misunder den Rige; Lykken fordærver de Lykkelige, Ulykken de Ulykkelige; Fordærvelsen
føder Forbittrelse, Forbittrelsen bryder
ud i Oprør.

Da Kadmus saaede Dragens Tænder, op
% --- p. 13 ---
\setcounter{footnote}{0}%
\opage{13}voxte harniskklædte Mænd, der dræbte hinanden:
disse harniskklædte Mænd ere Skjæbnens
Sønner.

Seet i Phænomenet, er det Menneskenes
Oprør imod hverandre indbyrdes; væsentlig
seet, er det den endelige, egoistiske Personligheds
Oprør mod Skjæbnen. Seet i Phænomenet,
er det en forandret Tilstand, en ny
Tingenes Orden, der fremgaaer af Kampen;
væsentlig seet, er det den samme Tilstand, den
samme forbigaaende, upaalidelige, falske Overeenskomst
mellem det Personlige og det Upersonlige, det Egoistiske og det Almene, altsaa
endnu som før: den samme Skjæbne.

I den borgerlige Spænding mellem det Selviske
og det Almene maa Personligheden strax
føle sig som den svagere; hvorfor den ogsaa
bevæger sig imellem de to Yderligheder: føielig
Underdanighed og fortvivlet Opsætsighed. Føler
Personligheden sig derimod ligesaa stærk som
det Almene, bliver Spændingen lidenskabelig
aandig, og Kampen \emph{heroisk}.

Den heroiske Villie nøies ikke med at modtage
Betingelser, den vil selv foreskrive Betingelser;
den tager ikke tiltakke med en taalelig
% --- p. 14 ---
\setcounter{footnote}{0}%
\opage{14}Tilstand, den fordrer ubetinget Tilfredsstillelse.
Her er desuden et høiere Mellemværende, hvori
de spændte Extremer, Personligheden og Skjæbnen,
forene sig, dette Mellemværende er Ideen.
Den heroiske Villie vil en Idee, og forsaavidt
det Almene; den vil tillige denne Idee virkeliggjort
i et bestemt verdsligt Maal, altsaa i en
Endelighed: her lurer Skjæbnen. Ved at afstikke
sig sit Maal maaler Personligheden sig
med Skjæbnen. At Skjæbnen har en Overmagt,
indlyser strax; men istedetfor at falde tilføie,
foreskriver Villien dog denne Overmagt
sine Betingelser, idet den foreskriver sig selv et
Valg. Dette Valg er imellem Alt og Intet,
Maalets Opnaaelse og Personlighedens Undergang.

„Det koste, hvad det vil, jeg gaaer til
Maalet!“ med dette Løsen har Friheden hævet
sig over Skjæbnen; saa tyder en Crøsus Oraklet,
saa læser en Wallenstein i Stjernerne, saa
gaaer en Cæsar over Rubicon, og saa begynder
Kampen.

Det er Maalets Opnaaelse, der giver det
Personlige sit Værd; heri ligger, „at, naar
Maalet tages bort, saa er Personligheden i sig
% --- p. 15 ---
\setcounter{footnote}{0}%
\opage{15}selv noget aldeles Ligegyldigt; den heroiske Beslutning
udtrykker altsaa den Modsigelse, at det
Personlige beviser sin høieste Gyldighed ved at
sætte sig som Middel for Maalet, som det i
sig selv Ligegyldige. Dersom denne Modsigelse
opgik for Bevidstheden i en klar og pludselig
Tanke, vilde den vel være som et Dolkestød,
der øieblikkelig dræbte Helten; men Modsigelsen
tilhylles af Ideen, fornemmes som Uro i Gemyttet, og forvandler sig til et Begivenhedernes
Drivhjul. Gjennemtrængt af Ideen, er den
personlige Selvfølelse fortabt i det Upersonlige.
Helten er allermeest sig selv, naar han er
allermeest udenfor sig selv i sin uimodstaaelige
Higen mod Maalet; greben af en uforklarlig
Magt, føler han sig mægtig; denne Magtfølelse
er hans Indvielse, hans mystiske Forbund med
Skjæbnen.

Det er Skjæbnens List at fange det Personlige
ved Personligheden selv, som det er
Skjæbnens Lyst at søndertræde sin Fange.
Skjæbnen bedaarer i Form af Lykke; den lammer
som den Gaade, der vækker Angst; den rammer
i Skikkelse af Nemesis; den bedaarede Crøsus,
den lammede Wallenstein, den knuste Cæsar.

% --- p. 16 ---
\setcounter{footnote}{0}%
\opage{16}Under den heroiske Kamp er det Personliges
Sammensmelten med det Upersonlige et Mysterium.
At afsløre dette Mysterium er en Opgave
for „den Vise“. Personligheden vender
sig da indad imod sig selv, for at klare og virkeliggjøre
sin frie Bestemmelse; saaledes bliver
Kampen moralsk.

For at maale sig med Skjæbnen maa Personligheden
allerførst gjøre sig Rede for, hvad
det da egentlig er, der staaer i vor Magt, og
hvad alt Det betyder, der ikke staaer i vor Magt.

Indsees det nu, at Fordelingen af de timelige
Goder og Onder: Sundhed og Sygdom,
Rigdom og Armod, borgerlig Frihed og borgerlig
Trældom o. s. v., ikke væsentlig staaer i
vor Magt, saa maa Frigjørelsen jo begynde
dermed, at Forskjellen imellem timelige Goder
og timelige Onder nedsættes til en i sig selv
ligegyldig Forskjel.

Det er altsaa „den Vise“ aldeles ligegyldigt,
om han er sund eller syg, rig eller fattig,
Herre eller Slave. Vistnok er det behageligere
at være sund end at være syg; men, da Villien
ligesaa godt kan bevise sin Kraft i Sygdom
som i Sundhed, kan Forskjellen mellem Be%
% --- p. 17 ---
\setcounter{footnote}{0}%
\opage{17}hageligt og Ubehageligt ikke røre det personlige
Selv.

At finde sig i det Ubehagelige og skikke sig
i Modgang er Resignation.

Allerede den borgerlige Kamp gjør Resignation
nødvendig. Den, der attraaer timelige
Goder, maa finde sig i timelige Tab; men det
er kun endelig Resignation, hvor Tabet af det
Ene erstattes ved Opnaaelsen af det Andet.
Saa vexler i Endelighed Frygt og Haab, Sorger
og Glæder, Tab og Erstatning; selv Tabet
af de Kjære, hvor smerteligt det end i Øieblikket
er, kan jo forvindes, erstattes --- med Tiden;
thi Tiden læger alle Saar.

Denne Endelighedens Trøst, denne timelige
Vexel af Frygt og Haab, af Øieblikkets Jammer
og Øieblikkets Jubel foragter „den Vise“.
Hævet i rolig, stolt Selvfølelse over det Omskiftelige,
resignerer han, ikke nu og da med
Hensyn paa dette eller hiint Endelige, men een
Gang for alle paa alt Endeligt, idet han erkjender,
at det Timelige efter sit Begreb er og
maa være timeligt, det Forkrænkelige forkrænkeligt.

„Har Du“ --- siger den Vise --- „et smukt
% --- p. 18 ---
\setcounter{footnote}{0}%
\opage{18}Leerkar, husk da vel paa, at det kan sønderbrydes.
Glemmer Du dette, da vil det smerteligt
forstyrre Dig, naar det skeer; men husker
Du altid derpaa, kan Tilfældet ikke overraske
Dig. Elsker Du din Hustru og din Søn, kom
da ihu, at de ere dødelige. Glemmer Du dette,
kan Smerten engang forvirre dit Sind; men
husker Du altid derpaa, da vil Du rolig, uden
Suk og Klage, følge dem til Graven.“

Hvad er saa Skjæbnen? I sine Phænomener
et vilkaarligt Spil af Medgang og Modgang,
Lykke og Ulykke, men efter sit Væsen det
i Tingenes Natur Begrundede, det Nødvendige,
det Almeenfornuftige. Heraf udledes den moralske
Lov, at den sædelige Personlighed ikke
maa ville Andet, end hvad Fornuften vil, og
ikke ønske Andet end det, der følger af Tingenes
Gang.

„Du skal ikke græde, naar Du sidder i
Theatret, og seer et Sørgespil; thi det er ikke
ørgeligt, at Helten falder, dersom det er nødvendigt.  % PRINTED AS IS: p. 18 „ørgeligt“ for sørgeligt (R12, read at the image)
Hvad der maa skee, det bør skee, og
dermed skal Du være tilfreds.“

Indviende sig til Nødvendighedens Tjeneste,
% --- p. 19 ---
\setcounter{footnote}{0}%
\opage{19}vil det personlige Selv forvandle Nødvendigheden
til Frihed og gjøre Friheden uimodstaaelig
som Nødvendigheden. Saaledes forsonet med
Tilværelsens Love, siger den Vise: „før mig,
o Skjæbne!“

For at virke i Verden maa Villien sætte sig
et verdsligt Maal; men her er Forskjellen mellem
„den Vise“ og „Daaren“. Daaren gjør
sin Personlighed til Middel for Maalet; den
Vise gjør ethvert endeligt Maal til Middel for
sin Personlighed.

„Verden“, siger den Vise, „er et stort Theater.
Menneskene ere Skuespillere. Guddommen
er den Digter, der lader Skuespillet opføre, og
fordeler Rollerne. Saa faaer den Ene Kongens,
den Anden Tiggerens, den Ene Herrens,
den Anden Slavens Rolle. Det tilkommer
ikke Skuespilleren at vrage sin Rolle, men
at spille sin Rolle. Kun Den, der spiller sin
Rolle godt, vinder Guddommens Bifald.“

Den Guddom, hvis Bifald den Vise attraaer,
er i egentlig Forstand den personificerede
Skjæbne; det moralske Selv er Skjæbnens
Personification. Den uendelige Resignation  er
% --- p. 20 ---
\setcounter{footnote}{0}%
\opage{20}altsaa dog en uendelig Egoisme; det er Skuespilleren,
der paà egen Haand har digtet sig
en Rolle; det er Skuespilleren, der spiller for
sig selv, og behager sig selv i Rollen; det er
Skuespilleren, der til Slutningen indhøster sit
eget Bifald.

At dette uden Overdrivelse forholder sig saaledes,
sees tydeligt i den sidste Akt.

Den skjæbnesvangre Personlighed, der her
har gjort sig til sin egen Guddom, kan ikke
tydeligere bevise sin Selvraadighed end derved,
at den raader over sit eget Liv, og derfor, selv
vælger og bestemmer sin egen Død. Det personlige
Liv er endeligt; det Endelige er ligegyldigt.
Med Frihed og Overlæg fuldbyrder
Skjæbnens Ven, hvad Skjæbnen ellers vilde
fremtvinge: han tager sig selv af Dage. Personligheden
hævder sig selv ved at tilintetgjøre
sig selv. Hermed ender den moralske Kamp og
det moralske Skuespil.

% --- p. 21 ---
\setcounter{footnote}{0}%
\opage{21}\afdeling{B.}{Om Forsynet.}
\stykke{}{Hvad er et Menneske?}

Ja, hvad er et Menneske? noget Stort og
Herligt; noget Ringe og Elendigt! Der er
paa Jorden ingen Skabning, der kan blive saa
lykkelig som et Menneske; thi der er paa Jorden ingen Skabning, der kan blive saa ulykkelig
som et Menneske; ingen, der kan glæde sig som
et Menneske, fordi ingen kan sørge som et
Menneske; ingen, der kan frygte, hvad et Menneske
kan frygte, fordi der er ingen, der kan
haabe, som et Menneske kan haabe; ingen, der
kan hade som et Menneske, fordi der er ingen,
der kan elske som et Menneske; ingen, der kan
hovmode sig som et Menneske; ingen, der kan
ydmyge sig som et Menneske; ingen, der kan
spotte som et Menneske; ingen, der kan tilbede
som et Menneske; ingen, ingen, der kan lide
den Uret, et Menneske kan lide; ingen, ingen,
der kan gjøre den Uret, et Menneske kan gjøre
mod sig selv og Andre. Hvad er da et Menneske?

% --- p. 22 ---
\setcounter{footnote}{0}%
\opage{22}Alle de Umyndige, der ere i Naturens Vold,
hjælpes af Naturen; Mennesket er i Skjæbnens
Vold: hvo hjælper Mennesket?

\delafsnit{Det aabenbarede Forsyn.}
\stykke{I.}{}
\motto{Thi han lader sin Sol opgaae over Onde og
Gode, og lader det regne over Retfærdige og
Uretfærdige.}

Men er det da en Hjælp? Kan da den
Gode trøste sig, at Solen opgaaer ogsaa over
de Onde? Kan den Retfærdige glæde sig, at
Himlen velsigner de Uretfærdige? Vi vide, at
det gaaer saaledes i Verden; vi vide, at Naturens
Love ikke kjende Forskjel paa Onde og
Gode, paa Retfærdige og Uretfærdige. Stormen
reiser sig, og Skibet synker, og Menneskene
forgaae; de forgaae baade Onde og Gode,
baade Retfærdige og Uretfærdige. Jorden aabner
sig, og opsluger Staden, og Menneskene
omkomme; de omkomme baade Onde og Gode,
baade Retfærdige og Uretfærdige. Pesten drager
igjennem Landene, og Menneskene døe af
den dræbende Sot; de døe baade Onde og
% --- p. 23 ---
\setcounter{footnote}{0}%
\opage{23}Gode, baade Retfærdige og Uretfærdige. Kan
Sligt da betyde et Forsyn? Det seer jo ud som
Skjæbne.

Ja, dersom alt Personligt er endeligt, og
det Evige selv upersonligt, saa er det Skjæbne.
Men, dersom det Personlige tillige er det Evige,
dersom den sande Personlighed er det Eviges
Sandhed: saa er der et Forsyn.

Dersom den evige Personlighed ikke selv var
menneskelig, var der for Mennesket intet Forsyn
og ingen Hjælp; dersom Mennesket af sig
selv var den evige Personlighed, da var hvert
Menneske sit eget Forsyn og i evig Forstand
sin egen Hjælper.

Hvert Menneske, der maa bekjende: „Jeg er
ikke Sandheden, men jeg trænger personlig til
Sandheden“, maa forholde sig troende til Forsynet,
og kan ikke være sit eget Forsyn.

Kun det Menneske, der turde sige om sig
selv: „Før Abraham blev, er jeg, det Eviges
Sandhed, den sande Personlighed“, kun et saadant
Menneske kan forholde sig vidende til Forsynet,
kan som menneskeligt Individ være under
Forsynet, og dog være Forsynet selv, der gjør
sig til Troens Gjenstand.

% --- p. 24 ---
\setcounter{footnote}{0}%
\opage{24}Den Menneskens Søn, der vidste med sig
selv, at han var det sande, i Kjød og Blod
aabenbarede Forsyn, og derfor gjorde sig til
Troens Gjenstand, har da heller ikke lært os
at see Forsynets Retfærdighed i en endelig
Fordeling af timelige Goder og timelige Onder.

At betale det Timelige med det Evige, det
Evige med det Timelige, er i al Evighed uretfærdigt;
men at betale det Timelige med det
Timelige, det Evige med det Evige, det er Lige
for Lige, det er guddommelig Retfærdighed.

Dersom det virkelig var Forsynets Tanke,
at det Gode og det Onde skulde gjengjældes
timeligt, maatte vi da ikke vente, at han, der
blev Menneske, aleneste for at han skulde være
det aabenbarede Forsyn, vilde faae baade Rigdom
og Magt og alle andre timelige Goder? Naar
det nu ikke skeer: hvad kan da være Grunden?
Skulde Forsynet maaskee misunde et Menneske
det Timelige? Eller skulde han, der er det
sande Menneske, sætte en Slags Ære i at undvære
og foragte det Menneskelige? Det er jo
en Spot over alt Guddommeligt og Menneskeligt.
Naar da han, der kunde have al Verdens
Herlighed, ikke har det, hvortil han kan
% --- p. 25 ---
\setcounter{footnote}{0}%
\opage{25}hælde sit Hoved, da er denne Fornedrelse ikke
at tilskrive en Villiens Vilkaarlighed, saa lidt
som en Skjæbnens Nødvendighed. Det er den
menneskelige Selvfornegtelse for Sandhedens
Skyld, det er Personlighedens uendelige Selvbekræftelse,
der i denne Ringhed aabenbarer den
guddommelige Overlegenhed, den fuldkomne Retfærdighed.

Alt i Aartusinder har Menneskeheden stridt
for timelig Gjengjældelse; og endnu gaaer Solen op over Onde og Gode, og endnu regner
det over Retfærdige og Uretfærdige. Den, der
verdslig seer paa Verdens Gang, seer i alt
Dette kun Skjæbnen; men Den, der troende
seer paa Menneskens Søn, seer Faderens Villie,
som er i Himmelen.

Dersom det var Forsynets Tanke, at Enhver,
der strider for en almeengyldig Idee,
skulde anerkjendes af Menneskene, og faae Ret
i Forhold til Ideens Værd og den Opoffrelsens
Kraft, hvormed han strider: maatte vi da ikke
vente, at Han, der ikke blot stred for den høieste
Idee, men var Ideen selv, skulde have tilkjæmpet
sig den høieste Anerkjendelse i Verden?

Alle de store Miskjendte, der ellers have
% --- p. 26 ---
\setcounter{footnote}{0}%
\opage{26}stridt for Ideer, og ere af Samtiden blevne
lønnede med Utak, Had og Foragt, kunde dog
kun fordre en betinget Anerkjendelse; det var
dog kun relative Ideer, for hvilke de strede; og
selv om En stred for det Høieste, kunde han
dog ikke sige med Sandhed: „Jeg, der strider
for det Høieste, er selv det Høieste.“

Menneskens Søn, der strider for den evige
Sandhed, vidner om sig selv: „Jeg er Sandheden“;
Menneskens Søn kan derfor gjøre sig
til Troens Gjenstand, kan fordre ubetinget
Underkastelse, guddommelig Ærefrygt.

Alle de store Miskjendte, der ellers kunne
klage over Menneskenes Utaknemmelighed, maae
dog vel tilstaae, at Samtiden ved at negte
dem Anerkjendelse i Ideens Navn, kun skader
sig selv til en vis Grad; Ideen kommer jo dog
alligevel engang Slægten tilgode: den Enkeltes
Personlighed er ligegyldig.

Menneskens Søn, der er Personligheden
selv, og maa fordre ubetinget Anerkjendelse for
Personlighedens Skyld, veed med uimodsigelig
og smertelig Vished, at hvert Menneske, der
negter ham, hvad han fordrer, tilføier sig selv
en evig personlig Uret.

% --- p. 27 ---
\setcounter{footnote}{0}%
\opage{27}Alle de store Miskjendte, der ellers fik Uret,
endskjøndt de havde Ret, kunde med Rette
klage over Skjæbnen; og vi, der nu indsee den
Uret, de maatte lide, vi kunne beklage dem, ja
vi kunne græde over dem, naar vi tænke paa
deres Skjæbne. Men Han, den Ene, der fik
Uret, ikke endskjøndt, men fordi han havde
Ret, evig, ubetinget Ret, han vil ikke beklages,
ikke ynkes, ikke begrædes; han siger jo udtrykkelig:
„Græder ikke over mig, men græder
over Eders egen og Eders Børns Skjæbne.“

Saa var da den Retfærdige iblandt de Uretfærdige;
Solen gik op over Onde og Gode, og
som Solen sank, gik Mørket op over Gode og
Onde; og Mørket styrkede de Onde, og i en
Mørkets Time led den Retfærdige for de
Uretfærdige. Det seer ud som en Skjæbne;
dog, det er ikke Skjæbne, det er Faderens Villie,
som er i Himmelen.

Dersom det var Forsynets Tanke, at Menneskene
selv skulde finde Veien ved Fornuftens
Lys, saa havde Menneskens Søn aldrig gjort
sig til Troens Gjenstand.

Kun det Lige erkjendes af det Lige, det
Personlige af det Personlige, men Guds og
% --- p. 28 ---
\setcounter{footnote}{0}%
\opage{28}Menneskenes Personligheder ere ikke hinandens
Lige.

„Agt din Ven, foragt din Fiende; skaf Dig
Ret i Tiden, om Du kan, thi i Evigheden holdes
der ingen Rettergang; Godt og Ondt, Ret
og Uret ere kun Endelighedens Differenser, personlige
Modsætninger, der alle opløses i det
evige, upersonlige Ene.“ Her er den Viisdom,
der negter Forsynet, og forguder Skjæbnen.

„Elsk din Ven, had din Fiende; thi dine
Venner ere den Eviges Venner, dine Fiender
den Eviges Fiender: see, Sværdet er i din
Haand, at Himmelens Retfærdighed maa fuldkommes
paa Jorden.“ Her er den Lidenskab,
der vil forvandle Forsynet selv til en
Skjæbne.

Personlighedens indbildte Overmagt er i
begge Former en evig Afmagt. Den, der foragter
sin Fiende, foragter væsentlig sig selv;
thi han foragter det Personlige: Foragtens
Overmagt er en Egoismens Afmagt. Den,
der gjør sin Fiende til Guds Fiende, raser
imod Gud: den fanatiske Overmagt er en vanvittig
Afmagt.

For at frelse den menneskelige Personlighed
% --- p. 29 ---
\setcounter{footnote}{0}%
\opage{29}fra dens naturlige Afmagt kom Forsynet selv i
Tjenerens Skikkelse. Var han kommen i Herskerens
Skikkelse med udvortes guddommelig
Overmagt, vare vi forblevne i vor naturlige
Afmagt. Men han fornegtede sig selv, lod
Verden beholde den udvortes Overmagt, og
blev lydig indtil Døden, at han i Selvfornegtelse
og Lydighed kunde aabenbare Personlighedens
uendelige Overlegenhed. Hans Fiender
vare visselig Guds Fiender; han saae deres
Forblindelse, han sørgede over deres Hjerters
Haardhed; men. han foragtede dem ikke, han  % PRINTED AS IS: p. 29 a small point after „men“ (seen at the image)
forbandede dem ikke: han bad for dem, og
døde.

I have hørt, at det er sagt: „Du skal
elske din Ven, og hade din Fiende“;
men jeg siger Eder: „elsker Eders Fiender,
velsigner dem, som forbande Eder;
gjører vel imod dem, som hade Eder,
og beder for dem, som krænke og forfølge Eder, paa det at I maae vorde
Eders Faders Børn, som er i Himlene.
\emph{Thi han lader sin Sol opgaae over
Onde og Gode, og lader det regne over
Retfærdige og Uretfærdige}.

% --- p. 30 ---
\setcounter{footnote}{0}%
\opage{30}\stykke{II.}{}
\motto{Ja endog alle Eders Hovedhaar ere talte, frygter
derfor ikke.}

Men ere disse Ord dog ikke at forstaae som
en digterisk Overdrivelse? Dersom den guddommelige
Retfærdighed virkelig er ophøiet over
enhver timelig Forskjel imellem Gode og Onde,
er det da tænkeligt, at den Evige skulde sænke
sig saa dybt i det Timelige og være saa nøieregnende
med det Enkelte? er det tænkeligt, at
Menneskehedens Forsyn skulde tælle alle Menneskenes
Børn, som en Hyrde tæller sine Faar?
er det den Alvidende værdigt at tælle saa smaaligt,
endogsaa Spurvene, endogsaa Haarene paa
Menneskets Hoved?

„Alle Eders Hovedhaar ere talte, frygter
derfor \textbf{ikke}!“ Dersom disse Ord kun vare menneskelige
Ord, da maatte vi vel betragte dem
som en Overdrivelse i et Begeistringens Udbrud;
thi i den Verden af Farer, der ligger
imellem Mennesket og Gud, har et Menneske
jo Meget at frygte. Men, naar Han, der
taler Ordet, er Ordet selv, da maae vi jo lære
af ham, hvor Faren er, og hvad vi have at
frygte.

% --- p. 31 ---
\setcounter{footnote}{0}%
\opage{31}Der er i Farernes Verden kun een ubetinget
Fare, kun Eet, der er frygteligt, kun Een,
der er at frygte. Derfor siger Ordet, det
aabenbarede Forsyn: „Jeg vil vise Eder, for
hvem I skulle frygte.“

Og Forsynet er menneskeligt. Han, der
taler om Faren, og viser os, for hvem vi skulle
frygte; Han, der overskuer alle Endelighedens
Farer og siger: „frygter ikke!“ han veed, hvad
Fare er; hans Liv var fortroligt med Faren.
Faren søgte ham i hans Vugge; den fulgte
ham til hans Grav. Overalt, i Ørkenen som
i Stæderne, paa Galilæas Sletter, imellem
Judæas Bjerge, overalt, hvor han gik hen, gik
Faren efter ham, og lurede paa ham.

Aarvaagen som Forsynet, frimodig som Den,
der er under Forsynet, gaaer Menneskens Søn
ad Farernes Vei, og frygter for Intet i Verden.
Han arbeider, og Intet standser ham;
han gaaer til sit Maal, og Ingen lægger Haand
paa ham, før hans Time kommer.

Men, da hans Time var kommen gik han
frivillig Faren imøde; den overraskede ham ikke
som et Tilfælde, den kom ikke over ham som
en Skjæbne; den kom, fordi han lod den komme.
% --- p. 32 ---
\setcounter{footnote}{0}%
\opage{32}Den kom da heller ikke i udvortes Larm, ikke i
en Storm af menneskelige Lidenskaber, ikke i et
Bulder af fiendtlige Vaaben; den kom i al
Stilhed, eensom til den Eensomme: kun saaledes
lod han den komme.

Ene hiin Nat i den stille Have, lod han
Faren komme til sig, og stred med den, ikke
som den Stærke strider, men som den Skrøbelige, at der i Kjød og Blod ikke skulde være
en Skjæbnens Forskjel imellem Stærke og
Skrøbelige, men en Hengivenhedens Seier i
Skrøbelighed.

Ene hiin Nat i den stille Have, lod han
Faren gaae som en Dødsangst igjennem sin
Sjæl; han trodsede den ikke, han foragtede den
ikke, han hævede ikke stolt sit Hoved som Den,
der offrer sig til Skjæbnen; han faldt paa sit
Ansigt, og bøiede sig ydmyg under Faderens
Villie.

Ene hiin Nat i den stille Have, lod han alle
Endelighedens Farer forvandle sig i Inderlighed
til een Fare. Saa stred han i Inderlighed,
og seirede i Inderlighed; og, da han saa reiste
sig fra Striden, see, da var det forbi med
Faren, da var der slet ingen Fare.

% --- p. 33 ---
\setcounter{footnote}{0}%
\opage{33}Dog --- hvad siger jeg? er det virkelig forbi
med Faren: nu kommer jo først den virkelige
Fare, thi nu kommer jo Vagten med Sværd og
Stænger. Ja, nu kommer Virkeligheden, men
--- Faren har vendt sig. Han, der spørger:
„Hvem lede I efter?“ er uforfærdet; Vagten
falder til Jorden af Forfærdelse.

Den Nat, da Menneskens Søn blev forraadt,
hudflættet og dømt, er i Virkeligheden
en Farernes Nat; men Faren har vendt sig.
Her er Fare for en Judas; her er Fare for
en Kajaphas. Her er Fare for en Pilatus,
der vadsker sine Hænder; her er Fare for
Folket, at den Uskyldiges Blod kommer over
dets Hoved. Ja, her er mange Farer og for
Mange; men for Ham, den Ene, for ham er
der ingen Fare.

Var der Fare for ham, da fik han vel Undsætning;
var hans Liv, hans Sag virkelig i
Fare, da skulde hans Tjenere vel stride for
ham, og mere end tolv Legioner Engle ruste
sig til hans Forsvar. Nu staaer han uden Forsvarere, vaabenløs, ene iblandt sine Fiender:
hvorfor? Fordi Faren alt er forbi; fordi hans
Hovedhaar ere talte.

% --- p. 34 ---
\setcounter{footnote}{0}%
\opage{34}„Ikke en Spurv falder til Jorden uden Faderens
Villie; ja endog Eders Hovedhaar ere
talte.“

Endelighedens Verden er en Tilfældighedernes
Verden. Fastholdes nu dette eller hiint
Tilfældige i dets tilfældige Enkelthed, da er
enhver saadan tilfældig Enkelthed, som f. Ex.
at en Spurv falder til Jorden, at der falder
et Haar af et Menneskes Hoved o. s. v., tilvisse
ligegyldig for Gud og Mennesker. En
saadan Fastholden af det Tilfældige i dets Tilfældighed
er nemlig aandløs, og det Aandløse
er i al Evighed ligegyldigt.

Det Tilfældige faaer først Betydning derved,
at det sees i væsentligt Forhold til et Andet,
som ikke er tilfældigt. Dette Andet kan
igjen være: enten Verdens Fornuftgrund eller
Personligheden selv.

Hvad Tilfældigheden betyder i dens Sammenhæng
med Verdens Grund og naturlige
Aarsager, er et videnskabeligt Spørgsmaal, der
aldeles ikke vedkommer Forsynstroen.

Hvad Tilfældigheden derimod betyder som
et verdsligt Mellemværende mellem Gud og
Mennesket, er et personligt Spørgsmaal. Dette
% --- p. 35 ---
\setcounter{footnote}{0}%
\opage{35}Spørgsmaal kan kun besvares af Ham., der  % PRINTED AS IS: p. 35 a point before the comma after „Ham“
aabenbarer Forsynets Tanke.

Her er en Blindfødt; lad os høre Mesterens
Samtale med Disciplene. „Og Disciplene
spurgte ham og sagde: Mester, hvo haver syndet?
denne eller hans Forældre, at han er født
blind? Og Mesteren sagde: hverken denne
syndede, ei heller hans Forældre; men paa det
at Guds Gjerninger skulde blive aabenbarede
paa ham.“

Disciplene og Mesteren ere altsaa enige i
den Forudsætning, at Tilfældet ikke kan være
noget Selvstændigt ligeoverfor Gud, at Forsynet
i ethvert Tilfælde hæver det Tilfældige.
Men, idet Disciplene ligefrem oversee Tilfældet,
hildes de i en dobbelt Misforstaaelse, forsaavidt
de paa een Gang miskjende Naturforholdenes
Realitet, og drage den guddommelige
Retfærdighed ned i det Timelige. Mesteren,
der kjender Forsynets Tanke, gjør opmærksom
paa Forskjellen imellem Naturbestemmelsen og
Guds udtrykkelige Villie. Forskjellen imellem
Seende og Blindfødte er en naturlig Forskjel i
det Timelige, der ikke ligefrem svarer til Guds
gjengjældende Retfærdighed. Forsaavidt ind%
% --- p. 36 ---
\setcounter{footnote}{0}%
\opage{36}rømmes altsaa det Tilfældige. Men det Tilfældige
er ikke for Gud i dets Tilfældighed;
det bestemmes som Middel for hans bestemte
Hensigt. Hvad der altsaa var Tilfældet med
hiin Blindfødte, at Guds Gjerninger paa ham
skulde blive aabenbarede, det Samme er just Tilfældet
med alle Blindfødte og med alle de mulige
Tilfælde i Verden, der vedrøre det Personlige.

--- „Medens jeg er i Verden, er jeg Verdens
Lys.“

Det naturlige Menneske, der seer Alt i Naturlighedens
Lys, seer Tilfælde paa Tilfælde,
saavel i Verdens Gang som i sit eget Liv; det
er det Blindfødte i den endelige Personlighed,
der ikke kan see Forsynet for lutter Tilfælde.
Menneskens Søn, den personlige Sandhed,
aabner den Blindes Øine, og hæver derved det
sørgelige Tilfælde.

„Medens jeg er i Verden, er jeg Verdens Lys.“

Forsynets Historie er en hellig Historie.
Dersom den hellige Historie var at opfatte ligefrem,
som enhver anden Historie, saa behøvede
vi kun at følge Menneskens Søn paa hans
Vandring, for at see, hvorledes han ophæver
% --- p. 37 ---
\setcounter{footnote}{0}%
\opage{37}Tilfældigheden, og aabenbarer Guds Gjerning.
Men i Forsynets hellige Historie er det Ydre
saa ueensartet med det Indre, at kun Den, der
seer Alt i Troens Lys, seer Forsynet.

Jo rigere Hjælpekilder Personligheden har i
sig selv, desto lettere forvandler den Tilfældet
til Middel og Anledning; jo mere overraskende
Tilfældet er, desto mere overraskende er ogsaa
den Aandsnærværelse, der gjennemlyner Tilfældigheden.
Derfor er Geniet i Pagt med
Lykken, er selv Lykken; derfor tæller den praktiske
Genialitet ligesaa mange Seire, som afgjørende
Tilfælde.

Dersom nu den personlige Overlegenhed,
der aabenbarer Forsynets Tanke, var at sammenligne
med Geniets Overlegenhed: da vilde
vi indbyde de Aandrige, de Geniale til at
gjennemgaae den hellige Historie; da vilde vi
sige til disse Aandrige, disse Geniale: „Seer
dog engang ret paa Menneskens Søn; som han
magter Tilfældigheden, har endnu Ingen paa
Jorden kunnet magte den; thi for ham er intet
Tilfældigt, han ophæver og forvandler alle
Tilfælde. Er det de mangehaande Sygdomstilfælde;
han hæver ethvert Tilfælde, thi han
% --- p. 38 ---
\setcounter{footnote}{0}%
\opage{38}helbreder jo alle de Syge. Er det de sørgelige
Dødstilfælde; han hæver dem, naar og
hvor han vil, thi han kan jo opvække de Døde.
Er det et af Ondskabens rænkefulde, spidsfindig
udtænkte Tilfælde, et Tilfælde, som med Skattens
Mynt, eller et Tilfælde, som med Qvinden
og hendes syv Mænd i Opstandelsen; see,
han magter alle disse Tilfældigheder. Sandhedens
Lyn splitter Ondskabens Væv, og Tilfældet
falder som en Skjæbne paa dem, der
forhærde sig i Usandhed.

Men, uagtet den personlige Magtfuldkommenhed
i alle disse Tilfælde er umiskjendelig,
er dog en saadan Opfattelse i sig selv en Bespottelse;
Menneskens Søn er ikke en Lykkens
Helt; han er Forsynets Hellige, Troens Gjenstand.
--- „Troer Du, at jeg kan gjøre dette?“ ---

Dersom Du ikke troer, da bliver Tilfældet
i dets sørgelige Tilfældighed; men „dersom
 Du kunde troe, da skulde Du see Guds
Herlighed.“
--- „Medens jeg er i Verden, er jeg Verdens
Lys.“

For Skjæbnens Blindfødte ere alle endelige
Tilfælde ligegyldige. Hermed opløses alt det
% --- p. 39 ---
\setcounter{footnote}{0}%
\opage{39}Enkelte i en eensformig, upersonlig, tom Almindelighed,
der mangler Livets Indhold og
Farvernes Rigdom. „Den Vise“ med al sin
Kraft er derfor ogsaa som et aandigt Skyggebillede,
et Skjæbnens Gjenfærd, der svinder i
det Dunkle.

For Den, der „troer paa Guds Søn“, den
personlige Sandhed, den sande Personlighed, er
det verdslig Tilfældige et personligt Mellemværende
mellem Gud og det enkelte Menneske.
Hvert mødende Tilfælde er en særegen Opgave,
Forsynet forelægger til personlig Løsning: saamange
Tilfælde, saamange Opgaver. Under
den endelige Løsning staaer Mulighed mod
Mulighed, Forstaaelse mod Misforstaaelse, Fare
mod Fare: --- Menneske, gjør hvad Du vil,
men gjør det i ubetinget Tro! Dersom Du
ikke troer, da bliver det Dig et sørgeligt Tilfælde,
men --- „dersom Du kunde troe, da
skulde Du see Guds Herlighed.“

--- „Medens jeg er i Verden, er jeg Verdens
Lys.“

Den udvortes vexlende Mangfoldighed af
endelige Tilfælde er, naar den sees i Sandhedens
Lys, lutter særegne Betingelser for Ud%
% --- p. 40 ---
\setcounter{footnote}{0}%
\opage{40}viklingen af det personlige Indre. Livets Tilfælde
ere Guds Tilskikkelser. Vantroen, der i
Tilskikkelsen kun seer Tilfældet, hjemfalder til
Skjæbnen; Troen, som i ethvert Tilfælde seer
Tilskikkelsen, frelses af Forsynet.

Menneskenes Forhold til det i den hellige
Historie sig aabenbarende Forsyn er derfor ikke
et eensformigt, farveløst Forhold, der taber sig
i upersonlig Almindelighed; det er et uendelig
indholdsrigt, mangfoldig, indtil de mindste Enkeltheder
nuanceret Livsforhold imellem den personlige
Sandhed og de menneskelige Personligheder.
Her staaer i hvert Optrin Tro mod
Vantro, de Seende mod de Blinde; her blive
de Blinde seende, de Seende blinde; her har
hver Begivenhed, hvert Tilfælde, hvert Ord,
hvert Blik, ja hvert Øieblik Uendelighedens
Betydning; her opgaaer Lyset over Onde og
Gode; her kommer Tilskikkelsen over Retfærdige
og Uretfærdige; her blive Guds Gjerninger
aabenbarede; her blive mange Hjerters
Tanker aabenbarede.

--- „Medens jeg er i Verden, er jeg Verdens
Lys.“

I den høieste Spænding mellem Viisdom%
% --- p. 41 ---
\setcounter{footnote}{0}%
\opage{41}mens Plan, Menneskenes Vilkaarlighed og Verdens
Tilfældigheder bliver Forsynet aabenbaret;
derfor er den hellige Historie en Lidelsernes
Historie.

Det seer i Lidelseshistorien ud, som vilde
Skjæbnen selv reise sig til Kamp mod Forsynet,
som skulde Tilfældigheden herske, og den menneskelige
Vilkaarlighed rive sig løs fra Styrelsen.

Og det seer i Lidelseshistorien ud, som havde
Guds Villie evig forudbestemt hver enkelt Begivenhed,
 hvert enkelt Optrin, hvert enkelt Træk.

--- Saa maa den Ene da „døe for Folket,
at ikke det ganske Folk skal fordærves“. Det
er et Ondskabens Indfald, en despotisk Vilkaarlighed
af en Enkelt i et enkelt Tilfælde; og
dog, --- det er jo ingen Vilkaarlighed, det er
jo en Aandens Indskydelse, det er jo en Spaadom,
der udtrykker Forsynets Tanke.
--- Saa bliver den Retfærdige da „regnet
iblandt Misdædere“; --- en oprørende, sørgelig
Vilkaarlighed, et sørgeligt Tilfælde! og dog,
det er jo ingen Vilkaarlighed; det er jo Viisdommens
forudbeskikkede Raad, det er jo Spaadommens
Opfyldelse.

% --- p. 42 ---
\setcounter{footnote}{0}%
\opage{42}Saa døer den Korsfæstede: Solen formørkes,
Forhænget revner, Gravene briste\dots{}.
„forunderlige Hændelser, maaskee opdigtede Tilfælde“\dots{}.
Ja, men --- dersom Du kunde
troe!

\emph{Sælges ikke to Spurve for een Penning,
og ikke een af dem falder til Jorden uden Eders Faders Villie. Ja
endog alle Eders Hovedhaar ere talte.
Frygter derfor ikke, I ere bedre end
mange Spurve}.

\stykke{III.}{For Gud ere alle Ting mulige.}

Men der skeer jo dog ingen Mirakler! ---
„Ja det er just det Betænkelige, at en Personlighed,
der tilhører Undernes Tidsalder, her
fremstilles som Forbillede, at Han, der gjør de
mange Undergjerninger, just skal være Exempel
for os, der leve under Naturlovene. Den, der
kan stille Stormen, kan jo gjerne sove paa
Skibet; men Den, der ikke kan stille Stormen,
maa vaage og bruge sine Kræfter. Den, der
med et Ord kan kaste Vagten til Jorden, kan
jo sagtens stikke Sværdet i Balgen; men, naar
% --- p. 43 ---
\setcounter{footnote}{0}%
\opage{43}Ordet Intet formaaer, maa man prøve det
blanke Vaaben.“

Dersom disse Betænkeligheder blot tilhørte
den erklærede Vantro, vilde vi slet ikke ændse
dem; men hvor er en Tro saa stærk, at den
ikke strider med Tvivlen?

De fleste Mennesker kjende vist saadanne
Sorgens, Ængstelsens, Spændingens Øieblikke,
i hvilke den Tanke ligesom uvilkaarlig paatrænger
sig: O, at der dog kunde skee et Mirakel!

--- „Han haver frelst Andre, sig selv kan
han ikke frelse.“ Saa lød det spottende Ord
over Ham, der havde hjulpet saa Mange, og
nu selv var som en Hjælpeløs.

Al Hjælp i det Timelige er i Forsynets
Tanke Middel for det Evige: der kan hjælpes
ved et Mirakel; men der kan ogsaa hjælpes
uden Mirakel. Bliver Miraklet forstaaet som
en Hjælp i det Timelige, saa hjælper det slet
ikke; det skader. Bliver Hjælpeløsheden i det
Timelige forstaaet som Vækkelsesmiddel for det
Evige, saa skader den slet ikke; den hjælper.
At der skeer et Mirakel, kan være baade til
Hjælp og til Fordærvelse; at Miraklet ikke skeer,
kan være baade til Fordærvelse og til Hjælp.
% --- p. 44 ---
\setcounter{footnote}{0}%
\opage{44}Et Menneske hjælpes kun af Gud, naar det
hjælpes ubetinget personlig; men det hjælpes
kun ubetinget personlig, naar det hjælpes til
Troen, i Troen, ved Troen. Hvorfor er
Miraklet skeet? At jeg skal troe. Hvorfor
skeer Miraklet ikke? At jeg skal troe. Den,
der negter Miraklets Mulighed, negter Forsynet;
Den, der antager, at Gud kun kan hjælpe
ved et Mirakel, negter ogsaa Forsynet.

Som Tegn paa sin guddommelige Magtfuldkommenhed
gjør Menneskens Søn Miraklet:
han bespiser de Hungrige. Men --- er det ikke
vidunderligt? --- Han, der hiælper Andre ved  % PRINTED AS IS: p. 44 „hiælper“ (the next line has hjælper)
et Mirakel, hjælper ikke sig selv ved et Mirakel.
Han, der formaaede at bespise fem Tusinde
i Ørkenen, forvandlede dog ikke Stene til Brød,
for at stille sin egen Hunger i Ørkenen.

Her sees det Dobbeltsidige i Forsynets Hjælp.
Den eensomme Hungrige, der fristes til et Mirakel,
og dog ikke vil hjælpe sig ved et Mirakel,
er dog alligevel væsentlig hjulpen: han har
personlig Hjælpen i sig selv; han har Hjælpen
i den guddommelige Vished, at et Menneske
ikke hjælpes paa een, men paa utallige Maader,
fordi „Mennesket ikke lever af Brød alene,
% --- p. 45 ---
\setcounter{footnote}{0}%
\opage{45}men af hvert Ord, som udgaaer af Guds
Mund.“

Den hungrige Mængde, der bliver hjulpen
ved Miraklet, er dog alligevel væsentlig hjælpeløs.
Den hungrer atter den næste Dag, og
søger Hjælperen, ikke fordi den saae Tegnet,
men fordi den aad af Brødet, og mættede sig.
Hjælpen blev tilbudt, og dog kunde det ikke
hjælpe; fordi de Hungrige troede, at det gjaldt
om Brødet alene, fordi Mængden troede, at
der skulde hjælpes timelig i det Timelige.

For Gud ere alle Ting mulige. Paa utallige
Maader kan Forsynet hjælpe; og dog er
det Altsammen forgjæves, naar Mennesket ikke
vil forstaae den evige, personlige Hjælp, og lade
sig hjælpe. Kunde der hjælpes timelig i det
Timelige, da skulde vel snart alle Hungrige mættes,
alle Nøgne klædes, og alle Endelighedens
Sorger endes; her ere jo fem Tusinde, som
kunne vidne, at Forsynet baade kan og vil
hjælpe; men her ere jo ogsaa fem Tusinde,
som kunne vidne, at dette ikke er Hjælpen.

For Gud ere alle Ting mulige. Paa overnaturlig
Maade kan Forsynet aabenbare sit
Herredømme over Liv og Død. Dersom de
% --- p. 46 ---
\setcounter{footnote}{0}%
\opage{46}Syges Helbredelse og de Dødes Opstandelse
var Hjælpen selv: da skulde Almagten vel snarlig
helbrede alle Syge, og aabne alle Grave; men,
naar det gjælder den forkrænkelige Personligheds
Forvandling til Uforkrænkelighed, kan der
kun hjælpes ved indvortes Helbredelse, indvortes
Opvækkelse. At Miraklet skeer, bliver i
Forsynets Aabenbaring netop et Beviis for, at
Miraklet i og for sig ikke er Hjælpen, og derved
tillige et Beviis for, at Hjælpemidlet ligesaa
godt kan bestaae deri, at der ikke skee Mirakler.

For Gud ere alle Ting mulige. Som Tegn
paa sin guddommelige Magtfuldkommenhed gjør
Menneskens Søn et Opvækkelsens Under. Men
--- er det dog ikke forunderligt? --- Han, der
veed med sig selv, at han kan vække den Døde
af Graven, see, han staaer som en Hjælpeløs
bedrøvet ved Graven og græder.

At kunne opvække den Døde og dog staae
og græde ved Graven, seer jo ud som en Modsigelse. Ja tilvisse, Ingen, der vil bruge sin
Forstand, skal kunne negte, at her maa være
en Modsigelse. Men, hvis er Skylden? hvo
er det, som her foraarsager Modsigelsen? Ikke
Hjælperen, men de mange Hjælpeløse, der ved
% --- p. 47 ---
\setcounter{footnote}{0}%
\opage{47}ham skulle hjælpes. Dersom de Hjælpeløse
strax kunde indsee, hvori Hjælpen bestaaer, da
var her Hjælp nok, da kunde Han, der staaer
ved Graven, i et Øieblik frelse Alle, og forherlige
sig selv. Men Hjælpen misforstaaes,
Midlet forvandles; hvad der burde være til
Hjælperens Forherligelse, bringer just Forfølgelsen
og Lidelsen over ham, gjør ham, menneskelig
talt, hjælpeløs; dette er just Modsigelsen,
og det en Modsigelse saa sørgelig, at
Guds Søn selv maa græde derover.

„Han haver frelst Andre, sig selv kan han
ikke frelse.“ --- Hvor der dog kan være Sandhed
i en Bespottelse! Thi det er jo bogstavelig
sandt, at han haver frelst Andre, og det er jo
ogsaa bogstavelig sandt, at han nu ikke kan
frelse sig selv --- ved Miraklet.

Havde Hjælpemidlet været Hjælpen selv;
havde Løgn og Ondskab, Vold og Uretfærdighed
kunnet tilintetgjøres ved et Mirakel: da var
Forsynets hellige Historie ikke bleven en Lidelsernes
Historie, men en Forherligelsens Historie
for Gud og Mennesker. Nu derimod, da Miraklet
kun er Middel, og Midlet forvendes,
saa hjælpes der ved det modsatte Middel, saa
% --- p. 48 ---
\setcounter{footnote}{0}%
\opage{48}hjælpes der i al Inderlighed derved, at der
ikke skeer noget Mirakel. For at hjælpe personlig
maa den personlige Hjælper selv blive
som den Hjælpeløse. ---“ „Og ville saa Alle blive
hjulpne?“

For Mennesker er det umuligt, men
for Gud ere alle Ting mulige.

\afdeling[\rule{10em}{0.4pt}]{C.}{Om Skjæbne og Forsyn.}
\stykke{I.}{Skjæbne er Skyld.}

Thi det, som man kan vide om Gud, er dem aabenbart, thi Gud har aabenbaret dem det, --- saa at
de ere uden al Undskyldning.

Den, der vil fortrøste sig til Forsynets
Hjælp, maa gjøre Forsynets Villie; Den, der
vil gjøre Forsynets Villie, maa kjende Forsynets Tanke. Hvo kjender Forsynets Tanke?
hvo gjør Forsynets Villie? Den personlige
Sandhed, den sande Personlighed.

% --- p. 49 ---
\setcounter{footnote}{0}%
\opage{49}Forholdet imellem den Ene, der er den personlige
Sandhed, og de Mange, der ved ham
skulle deelagtiggjøres i Sandheden, er ikke et
udvortes historisk Forhold imellem et enkeltstaaende
Individ, der engang har været til, og
en Mængde andre Individer, der i Tiden have
sluttet sig sammen om denne Ene. Var det saaledes,
da vilde der jo være en Skjæbnens Forskjel
imellem dem, der vare saa lykkelige at
komme i personligt Forhold til det overordentlige
Individ, og de Mange, Mange, der vare
lovlig undskyldte, at de ikke kjendte ham, fordi
de aldrig havde seet ham, og aldrig hørt et
Ord om ham.

Den personlige Sandhed er Sandhedens
Konge; hver Den, som elsker Sandheden, hører
hans Røst, og hver Den, som hører hans Nøst,  % PRINTED AS IS: p. 49 „Nøst“ for Røst
frelses af Forsynet.

Al Skjæbne er personlig Usandhed; thi Skjæbnen
kommer kun derved, at det Personlige selv
gjør sig til det Usande. Al Skjæbne er væsentlig
Skyld; thi Skjæbnen kommer kun derved, at det Personlige krænker sig selv ved at
bøie sig under det Upersonlige. Al Skjæbne er
væsentlig Synd; thi Skjæbnen opkommer kun
% --- p. 50 ---
\setcounter{footnote}{0}%
\opage{50}derved, at Mennesket fornegter Gud ved at
forvandle ham til et Upersonligt.

Personlighedens umiddelbare Selvfornemmelse
bestemmes deels ved det Endelige, deels
ved det Evige. Selvfornemmelsen, umiddelbart
bestemt ved det Endelige, er Egoisme; bestemt
ved det Evige, er den Samvittighed.

Egoismens Ideal er en Blanding af det
Personlige med det Upersonlige; denne Blanding
er det Forblommede, hvori den endelige
Personlighed skjuler sin Afmagt for sig selv;
dette Forblommede er det Mythiske. Det er
ikke blot Afguden, der er mythisk; det Mythiske
er i Afgudsdyrkerens upersonlige Personlighed.\footnote{Det \textit{Mythiske} er ikke en Phantasiens Bestemmelse; det er en ethisk Bestemmelse. Den mythiske Fortolkning saavel af det Gl. som af det N. Tst. er begrundet i personlig Usandhed, eller, om man vil, i upersonlig Sandhed.}

Samvittighedens Ideal er en Eenhed af Aabenbaring
og Tro. I Aabenbaringen bliver den
egoistiske Afmagt sig selv aabenbar ligeoverfor
den personlige Almagt, og i Troen faaer Mennesket
sin sande Personligheds Kraft som en
Guds Kraft.

% --- p. 51 ---
\setcounter{footnote}{0}%
\opage{51}Alle Egoismens Idealer staae under Skjæbnen;
alle Samvittighedens Idealer ere under
Forsynet.

Igjennem hele Historien gaaer Striden imellem
personlig Sandhed og personlig Usandhed,
imellem Egoismens og Samvittighedens Idealer. Under Stridens Gang skiller Slægten sig
ad, og Bevægelsen deler sig i tvende Strømme:
den mythiske og den i Aabenbaringstroen \emph{religiøse}.

I den mythiske Retning seirer Egoismens
Idealer under idelige Modvirkninger af Samvittigheden;
i den troende Retning seirer Samvittigheden
under idelige Modvirkninger af
Egoismen. Adskillelsen er gjennemgribende som
en Adskillelse imellem tvende, hinanden modsatte
Aandssphærer. Over begge er der en Skjæbne,
thi i begge er der en Skyld; men ogsaa over
begge er der et sig forbarmende Forsyn, der
straffer Skylden.

For Mythernes Folk er Straffen Skjæbne;
for Aabenbaringens Folk er Skjæbnen Straf;
begge føres til Maalet: den personlige Sandhed
er Forsynets Maal, og Maalet naaes i
Tidens Fylde. 

% --- p. 52 ---
\setcounter{footnote}{0}%
\opage{52}I Tidens Fylde optræder altsaa det virkeliggjorte
Samvittighedsideal: \emph{frelsende}, thi
Personlighedens Frelse er Samvittighedens Sag;
\emph{dømmende}, thi hvert Menneske skal dømmes
af sit Ideal. Men Idealets Sandhed er Sandhedens
Konge; hvert Glimt af personlig Sandhed,
hver Gnist af Samvittighed, der yttrer sig
i Verden, er en Gnist af hans Flamme, et
Glimt af hans Lys.

Den personlige Sandhed er Sandhedens
Konge; hver Den, som elsker Sandheden, hører
hans Røst, og erkjender, at han er Idealet selv.
Men de, som forblive i Usandheden, høre ikke
hans Røst; de ville helst beholde deres egne
selvgjorte, forqvaklede Idealer, den Ene i Naturen,
den Anden i den skrevne Lov. Nu bliver
Skylden aabenbar. Brøden ligger nu ikke
deri,. at Idealet er blevet formørket og forvansket  % PRINTED AS IS: p. 52 a point after the comma
i Tidernes Løb; thi denne Forvanskning
hæves jo netop derved, at Sandhedens Konge
selv viser sig i Verden; nei, Brøden ligger just
deri, at enhver af Parterne paa sin Viis vil
beholde det forqvaklede Ideal, det Ideal, der
modsiger Samvittigheden, men smigrer Egoismen.
Saaledes bliver da Miskjendelsen af
% --- p. 53 ---
\setcounter{footnote}{0}%
\opage{53}Forsynet til aabenbart Oprør imod Forsynet;
et saadant Oprør er en Forblindelse, Forblindelsen
er Skjæbne, Skjæbnen er Skyld: her
gjælder ingen Undskyldning.

Sættes Sandheden upersonlig, da er det
ganske rigtigt, hvad den moderne Aand paastaaer,
at Menneskeslægtens Historie er syndfri,
og at Synderne falde paa alle Individer
uden Undtagelse; men anerkjendes den personlige
Sandhed, saa gjælder Troens Ord, at
Slægtens Historie er en syndig Historie, og at
der er Een i Slægten, der er syndfri.

Verdenshistorien er i uophørligt Oprør mod
Forsynet; men Styrelsen viser sig deri, at
Oprørerne ere afmægtige. For at see Oprøret
i hele dets Udstrækning maae vi skjelne mellem
de larmende Oprørere og de stille Oprørere.

I et Aandens Billede skildres de larmende
Oprørere fra Oldtids Dage som Folkene, „der
bruse og pønse paa tomme Raad, som Jordens
Mægtige, der ville sønderrive Styrelsens
Baand og stride mod Forsynets Konge“. Mod
disse bruges det stærke Ord, at Han, der
sidder i Himmelen, leer.

Straffen, der skal ramme disse Oprørere,
% --- p. 54 ---
\setcounter{footnote}{0}%
\opage{54}kommer da ogsaa i Skjæbnens Skikkelse; thi
Forsynets Konge skal „knuse dem med Jernstav,
som en Pottemagers Kar skal han slaae dem i
Stykker.“

De larmende Oprørere ere udvortes kjendelige,
just fordi de larme; anderledes forholder
det sig med de stille. De stille Oprørere reise
ikke dristig Ideal mod Ideal; de hæve ikke
Hovedet saa høit, at Skjæbnen historisk kan
ramme dem: hvilke ere da de stille Oprørere?

Lad os først tilstaae, at det stille Oprør ulmer
i os Alle, og derhos charakterisere Oprørerne
selv efter Begrebet. Det er jo alle de, der
fortrøste sig til Forsynets Bistand uden nogensinde
at forelægge sig selv det alvorlige Spørgsmaal: hvad er da egentlig Forsynets Tanke?
det er jo alle de, der trøstig beraabe sig paa
Sandhed og Ret, uden dog at ville bøie sig
ubetinget for Retfærdighedens Ideal, og høre
Sandhedens Konge; det er jo alle de Mennesker,
Høie og Lave, Rige og Fattige, Begavede
og Ubegavede, der ene og alene, eller dog nærmest
og allernærmest, indrette deres Liv paa
hyggeligt Velvære, behagelig Tryghed, endelig
Travlhed, timelig Virksomhed for timelig For%
% --- p. 55 ---
\setcounter{footnote}{0}%
\opage{55}maal, og for Resten alle ere dydzirede, forsaavidt
de alle ere zirede med de for den jordiske
Velværen høist nødvendige Dyder, saasom:
Skikkelighed, Ærbarhed, Vindskibelighed,
Ædruelighed, Redelighed i Pengesager, Fredsommelighed,
Tjenstvillighed, hvortil endnu kan
regnes Godmodighed, Omgjængelighed, Selskabelighed
o. o. v., alle disse, der have alle
disse Dyder og hverken meer eller mindre, alle
disse ere hverken meer eller mindre end honette,
ærbare, stille, skikkelige Oprørere imod den evige
Forsynets Tanke, imod den Sandhedens Konge,
der fornedrede sig selv, og lod Tilskikkelserne
komme over sig, at det dog maatte blive aabenbart,
hvad det er, at kjende Forsynets Tanke
og gjøre Forsynets Villie.

\stykke{II.}{Timelige Vilkaar.}

Herren er den, som gjør fattig, og den, som gjør rig;
han er den, som nedtrykker, og den, som ophøier.

Uligheden i de menneskelige Vilkaar er vistnok
i verdslig Forstand at betragte som et sørgeligt
Misforhold, en Skjæbnens Uretfærdighed;
men i Troens Forstand er denne Ulighed
% --- p. 56 ---
\setcounter{footnote}{0}%
\opage{56}just et Beviis paa Forsynets allestedsnærværende
Retfærdighed.

I timelig Henseende har den Rige. en overordentlig  % PRINTED AS IS: p. 56 full stop before lower-case „en“ (R04)
Begunstigelse fremfor den Fattige, den
Mægtige fremfor den Ringe; i Verdslighedens
Sprog maae vi udtrykke dette saaledes: Skjæbnen er den, som gjør fattig, og den,
som gjør rig; det er den, som \emph{nedtrykker}, og den, som ophøier.

See vi derimod Sagen i Troens Lys, da
er heri slet ingen Uretfærdighed, ingen Vilkaarlighed,
ingen Partiskhed. At dette ikke er
en vilkaarlig Paastand, men en evig Sandhed,
følger ligefrem af Modsætningen mellem det
Verdslige og det Religiøse.

Ligeover for Skjæbnen, altsaa i Verdsligheden,
er Menneskets Betydning i enhver Henseende
endeliggjort. Et Menneske kan i Verden være noget Stort; men al verdslig Storhed
er relativ: selv den største Individualitet
forsvinder som en Ubetydelighed i det umaadelige
Hele. Omvendt kan et Menneske være
noget saare Ringe; men selv den Ringeste er
dog altid Noget. Som nu al Storhed og al
Ringhed i verdslig Forstand er Stykværk, saa er
% --- p. 57 ---
\setcounter{footnote}{0}%
\opage{57}ogsaa al Menneskets Betydning Stykværk: Verdsligheden
er just det udstykkede Menneskevæsen.

Det er ikke den Samme, der paa een Gang
er baade stor og ringe; men den Ene er stor,
den Anden ringe; det er ikke den Samme, der
paa een Gang er baade rig og fattig; men den
Ene er rig, den Anden fattig; ja den Samme
er end ikke altid det Samme: før var han
mægtig, nu er han kun ringe; før var han
fattig, nu er han jo rig.

Fastholdes nu denne Menneskevæsnets verdslige
Udstykning i dens Verdslighed, saa have
vi Skjæbnen. Hvert Menneske, der har sit
Liv i denne Udstykning, udtrykker personlig den
upersonlige Sandhed: Skjæbnen er den, som
gjør fattig, og den, som gjør rig; det er den,
som nedtrykker, og den, som ophøier.

Ligeover for Gud er Mennesket derimod i
dets Endelighed „Støv og Aske“, en ligegyldig
Størrelse, et reent Intet; og dog er Mennesket
i Guds Tanke noget uendelig Stort, Aand af
Guds Aand, Guds Billede. Her falder til
begge Sider det Endelige, det Relative bort,
og dermed falder ogsaa Udstykningen bort.

Her er den Samme væsentlig altid det
% --- p. 58 ---
\setcounter{footnote}{0}%
\opage{58}Samme. Den Samme, der her er ringe, er
ogsaa stor, thi hans Storhed fremgaaer af
hans Ringhed; den Samme, der her er fattig,
er ogsaa rig, thi i hans Fattigdom fødes personlig
en Evighedens Rigdom.

Men nu Stykværket, Verden og Verdsligheden?
Verdsligheden maa bort; men Verden
bliver, thi Verden er god; Stykværket bliver,
thi dette Stykværk er saare godt --- i Forsynets
Øine.

Herren er den, som gjør fattig, og
den, som gjør rig; han er den, som
\emph{nedtrykker}, og den, som ophøier.

Alle Tilskikkelser ere personlige Opgaver. Formedelst
Stykværket bliver den ene Opgave til
mange Opgaver, og just i denne Forskjellighed
for hvert Menneske en virkelig sand og personlig
Opgave.

Hvad der i hver Tilskikkelse er guddommelig
Hjælp, er tillige timelig Byrde, og Alt,
hvad der seer ud som timelig Hjælp, er omvendt
guddommelig Byrde.

Stille vi den Mægtige og den Ringe imod
hinanden, da kommer Opgaven i aldeles modsatte
Tilskikkelser, og, saa at sige, under to For%
% --- p. 59 ---
\setcounter{footnote}{0}%
\opage{59}klædninger. Enhver af de Tvende faaer umiddelbart
det personlige Indtryk af Tilværelsen,
der svarer til hans Stilling: et Menneske er
noget Stort; et Menneske er noget Ringe.

Dersom nu Opgavens Løsning blot bestod
i en Tilegnelse af det personlige Indtryk, at et
Menneske er noget Stort, og kan udrette noget
Stort i Verden: da var den Mægtige vilkaarlig
begunstiget, thi det Indtryk modtager han jo
netop fra alle Sider i hvert Øieblik, han lever.

Dersom Opgavens Løsning omvendt bestod
i en umiddelbar Tilegnelse af det Indtryk, at
et Menneske er noget Ringe; da var den Ringe
jo ligefrem begunstiget ved sin Ringhed, thi
Verden skal nok lære ham, hvor ringe han er.

Men --- her er den Samme væsentlig altid
det Samme. Opgavens Løsning bestaaer altsaa for begge i en Forarbeidelse og inderlig
Tilegnelse af begge Indtryk; herved blive Vilkaarene
lige.

Kan nu den Mægtige forarbeide Indtrykket
af verdslig Storhed saaledes, at han --- ikke
blot til Opbyggelse i en stille Time, men i al
verdslig Omgivelse, i virkelig Personlighed --- i
hvert Øieblik forvandler dette Indtryk i religiøs
% --- p. 60 ---
\setcounter{footnote}{0}%
\opage{60}Følelse af uendelig Ringhed for Gud; kan den
Mægtige, naar Tusinder med Beundring see op til
ham, dog endnu altid see det Øie, der prøvende
seer ned paa ham: da løser han en stor Opgave;
da forstaaer han, at timelig Hjælp er guddommelig
Byrde; da vil han ogsaa bære sin Byrde,
blive i Høiheden, og opløfte sit Hoved som Den,
der er sat i Høiheden, forklarende den Tanke,
at Herren er den, som ophøier.

Saa kommer den Ringe, men --- hvor stor,
hvor ophøiet er ikke hans Opgave! Hvad han
er, hvad han føler, tænker og gjør, er jo dog
saa ubetydeligt, for det Hele saa ligegyldigt,
at Ingen lægger Mærke dertil, fordi det slet
ikke er værd at bemærke. Og dog skal denne
Ubetydelige og alt det Ubetydelige, han gjør,
være af uendelig Betydning. For ham, der end
ikke er kjendt i Verden, skal hvert Øieblik være
ligesaa vigtigt, ligesaa spændende, ligesaa afgjørende,
som for Den, der er stillet saa høit,
at hans Vink bevæger en Verden. Kan nu
den Ringe løse denne Opgave, kan han, den
Ubemærkede, virkelig tilegne sig det personlige
Indtryk, at han i ethvert Nu bemærkes af Forsynet: da vil han i Sandhed forstaae, at Det,
% --- p. 61 ---
\setcounter{footnote}{0}%
\opage{61}der er timelig Byrde, er guddommelig Hjælp;
da vil han, naar det saa skal være, gjerne blive
i Ringheden som den Fortrykte: han trykkes jo
ikke af Skjæbnen; han udtrykker jo den Tanke,
at Herren er den, som nedtrykker.

Sammenstille vi nu den Mægtige med den
Ringe, saaledes som de selv begge have stillet
sig under Forsynet; da vil den Første vist ikke
foragte den Sidste, og den Sidste ikke misunde
den Første; nei, nu ville de forstaae hinanden
i al Inderlighed, og hjælpe hinanden af al Kraft,
den Mægtige ved sin Magt, den Ringe ved
sin Ringhed.

„Herren er den, som gjør fattig, og den,
som gjør rig.“

Naar de Rige i taknemmelig Erkjendelse af
Forsynets Viisdom indskjærpe denne Sætning
som guddommelig Garanti for Eiendomsretten,
medens de Fattige i Utaalmodighed over Skjæbnens
Vilkaarlighed reise sig, for at komme Viisdommen
til Hjælp: saa bruse Vandene, saa
mødes Oprørerne fra begge Sider, de stille med
de larmende; men --- Han, der sidder i Himmelen,
leer.

Det er i den almindelige Mening en let
% --- p. 62 ---
\setcounter{footnote}{0}%
\opage{62}Sag at takke Forsynet for Rigdom, men saare
tungt at skulle takke for Fattigdom. Dette
er imidlertid ikke den guddommelige Mening.
\emph{Sandheden} siger: „salige ere de Fattige“;
Sandheden siger: „det er lettere, at en Kameel
gaaer igjennem et Naaleøie, end at en
Rig kommer i Guds Rige“. Dersom det da
er sandt, hvad \emph{Sandheden} siger, saa maa
Rigdom jo være en stor og forfærdelig Byrde;
men dersom Rigdom er en saa stor og forfærdelig
Byrde, da er den Rige jo ikke af Gud
begunstiget fremfor den Fattige.

Hvad skal den Fattige gjøre? Forstaae Guds
Tilskikkelse, og forvandle Byrden til Hjælp.

Hvad skal den Rige gjøre? Forstaae Tilskikkelsen,
og forvandle Byrden til Hjælp. Men
saa skal jo den Rige gjøre det Selvsamme som
den Fattige. Skulle de nu begge, hver paa sin
Maade, gjøre det Samme; da trænge de ogsaa til gjensidig Understøttelse, og bør følgelig
hjælpe hinanden; men især bør den Fattige
hjælpe den Rige, thi den Fattige har jo det
Mindre at bære.

% --- p. 63 ---
\setcounter{footnote}{0}%
\opage{63}Den Fattige, der hjælpes af Gud saaledes,
at han --- i personlig Sandhed virkelig hjælper
den Rige, beviser just derved i Gjerning, at
Herren er den, som gjør fattig.

„Det er lettere, at en Kameel gaaer igjennem
et Naaleøie, end at en Rig forstaaer Forsynets
Tanke og gjør Forsynets Villie“; Rigdommen
gjør altsaa Opgavens Løsning saa vanskelig, at
Byrden næsten maa betragtes som overmenneskelig.
Hvad er da Sandhedens personlige Raad
til den Rige? At den Rige snarest muligt formindsker
sin overmenneskelige Byrde, og hjælper
sig ved at hjælpe den Fattige. Kun naar den
Rige saaledes hjælper den Fattige, kun naar
den Rige personlig for Gud erkjender, at han,
hver Gang han afhjælper den Trængendes
Nød, væsentlig hjælper sig selv, udtrykker han
den Sandhed, at Herren er den, som gjør
rig.

Saa ere da Vilkaarene lige; saa ere vi da,
skjøndt under mange Forklædninger, alle een
Faders Børn; saa er da Retfærdigheden guddommelig,
som Gud er retfærdig.

% --- p. 64 ---
\setcounter{footnote}{0}%
\opage{64}\stykke{III.}{Aandelige Vilkaar.}

Al god og fuldkommen Gave kommer herovenfra fra
Lysenes Fader, hos hvem der ikke er Forandring
eller Skygge af Omskiftelse.

Vil Nogen spørge ligefrem, om Talenterne,
Genierne, Heroerne og de store Verdensideer
dog ikke ifølge deres høiere Natur maae siges i
ganske særegen Forstand at være under Forsynet;
da kan et saadant Spørgsmaal ikke besvares,
saalænge det holdes i en Almindelighed,
der udelukker det Personlige. Først, naar det
Personlige tages med, bliver der Mening i
Spørgsmaalet, og Besvarelsen vil da beroe
paa den Stilling, Personligheden indtager til
Ideerne.

Den, der holder sig til den upersonlige
Sandhed, altsaa i dette Tilfælde --- hvor det
Personlige netop er det Afgjørende --- til den
personlige Usandhed, vil da naturligviis paastaae,
at Ideerne oprinde af Verdensfornuftens
Naturgrund, og at de høiere Aandsgaver følgelig ikke ere Andet end rene Naturgaver.

Den derimod, der troer personlig paa den
personlige Sandhed, vil ubetinget fastholde den
% --- p. 65 ---
\setcounter{footnote}{0}%
\opage{65}Forudsætning, at al god og fuldkommen Gave
kommer herovenfra fra Lysenes Fader, at altsaa Ideerne i Historien saavelsom Fuglene under
Himmelen, Genierne i Kunsten saavelsom
Lilierne paa Marken staae under Guds faderlige
Varetægt.

Forskjellen imellem Begavede og Ubegavede
er i en vis Maade at sammenligne med Forskjellen
mellem Rige og Fattige. De Begavede
ere jo de aandelig Rige, der meddele os af
deres Overflødighed, hvad de godt kunne undvære.
Uligheden faaer imidlertid en ganske anden Betydning i det Aandelige end i det Timelige.
Seet i Ideen, er Intet saa agtværdigt,
som et sandt Aandsaristokrati, og seet i Ideen,
er Intet saa foragteligt, som det rene Pengearistokrati.

De Aandsbegavede ere bestemte til at være
Slægtens Velgjørere, Lærere, Opdragere og
Tugtemestere; de ere Ideens Tjenere, udsendte
til Tjeneste for dem, der skulle deelagtiggjøres
i Ideernes Rigdom; men hvilken af disse
Tjenere, vi end vælge os til Exempel, saa
overbevises vi om, at enhver overordentlig
Gave forvandler sig til en guddommelig Byrde,
% --- p. 66 ---
\setcounter{footnote}{0}%
\opage{66}naar den Begavede skal udtrykke Forsynets
Tanke.

For Gud er ethvert Menneskeværk, selv det
største Mesterværk, Intet. For Talentet er
Værket Alt; thi det er jo „Værket, der skal prise
Mesteren“. For Gud gjælder alene Samvittighedens
personlige Ideal; i de relative Aandssphærer
gjælde de mange upersonlige, til Værkernes
Udførelse ledende Idealer. Ville de
Begavede nu stille sig under Forsynet, da vil
jo Enhver af dem komme til at leve for
tvende, hinanden modvirkende, Idealer; thi Enhver
af dem faaer da foruden Samvittighedsidealet,
der gjør Værket til et Intet, endnu
desuden det til hans Talent svarende Ideal,
der gjør Værket til Alt. Men at leve for
tvende, hinanden modstræbende, Idealer er overordentlig byrdefuldt.

At de Begavedes Fortrin fremfor de Ubegavede
virkelig er en guddommelig Byrde, kan
paavises i Exempler fra de forskjellige Aandssphærer.

Lad engang den geniale Kunstner stille sig
under Forsynet ved Siden af den simple Arbeidsmand. De anerkjende da begge Samvit%
% --- p. 67 ---
\setcounter{footnote}{0}%
\opage{67}tighedsidealet; de ere altsaa lige i religiøs Begeistring.
Men Arbeidsmandens Værk ligger
saa langt under det Ideale, at det umulig
kan blive Gjenstand for fristende Begeistring;
Kunstnerens Arbeide er derimod umuligt uden
en tilsvarende Begeistring.

Hvilken Forskjel! Naar Arbeidsmanden kun
dyrker een Gud, og kun har een Begeistring,
faaer Kunstneren jo næsten to Guder, idet han
faaer to Idealer, og begeistres i to ueensartede
Retninger. Medens der ikke let skal findes
nogen Arbeidsmand, der mener, at hans Værk
som Værk er til Guds Ære (thi det vilde dog
være besynderligt, om En paastod, at han kunde
grave Gruus til Guds Ære, eller slaae Kalk til
Guds Ære), saa er Kunstneren idelig fristet til
at sammenblande Kunstidealets og Samvittighedsidealets
Fordring. Ja, dersom det blot
kunde lykkes Talentet at forvandle Kunstdyrkelsen
til en ligefrem Gudsdyrkelse, saa var
Kunstneren hjulpen, saa kunde man rigtignok
sige, at den Begavede var baade Lykkens og
Guds Yndling. Men det er umuligt: det
æsthetiske Pietisteri er Oprør mod Forsynet.
Der kan ligesaalidt males et Billede,  være sig
% --- p. 68 ---
\setcounter{footnote}{0}%
\opage{68}et catholsk Mariebillede eller et protestantisk
Christusbillede, til Guds Ære, som der kan
componeres en Ouverture til Guds Ære, eller
spilles Comoedie til Guds Ære.

Al religiøs Kunstdyrkelse er raffineret Afgudsdyrkelse.
I Kunsten er Værket Alt; i Forsynets
Øine er Værket Intet. Adskillelsen af
hellig og verdslig Kunst er saa langt fra at
betegne en virkelig Forskjel imellem det Religiøse
og det Profane i personlig Forstand, at
den „hellige“ Kunst netop medfører den største
Fristelse til at profanere det Hellige. Forskjellen er vel hentet fra Stoffet, men dog fornemmelig
fra Talentets Beaandelse. Den profane
Kunstner (Skuespildigteren, Skuespilleren) beaandes
verdslig, og er ikke saa let fristet til at
forvexle sin Muse med en Forsynets Engel;
den „hellige“ Kunstner faaer derimod sin Beaandelse
gjennem de hellige Gjenstande, hans
Kunstbegeistring er som en Gudsbegeistring:
saa skeer det let, at den Begeistrede bliver
svimmel, annammer Indgivelsen som en Guds
Aabenbaring, falder tilbedende paa sine Knæ,
og bliver --- en raffineret Afgudsdyrker.

Ethvert overordentligt Talent er en god og
% --- p. 69 ---
\setcounter{footnote}{0}%
\opage{69}fuldkommen Gave; men det koster utrolig Anstrengelse
at udtrykke den personlige Sandhed,
at Gaven kommer herovenfra.

Den Ikke-Begavede, der æder sit Brød i sit
Ansigts Sved, og sukker under Endelighedens
Byrder, har i disse Byrder en guddommelig
Hjælp, og fristes ikke af de relative Idealer.
Den Begavede, der bevæges i Stemning og hengiver
sig i ideel Productivitet, maa fordobble sin
Anstrengelse for at bevare sin Personlighed.
Værket, der er ham Alt, maa tillige være ham
et Intet, et Ligegyldigt, en Barneleg, et Vehikel
for Udvikling af ethiske Kræfter. Den
Begavede, der nu kan gjøre dette: kan have
Kunsten, som om han ikke havde den, kan elske
dens Ideal, som om han ikke elskede det, fordi
han elsker Gud over alle Ting, kan takke Gud
ligesaa inderlig, naar Værket aldeles mislykkes,
som naar det over al Maade lykkes, kan holde
sig saa uforandret til Samvittighedens Ideal,
at der i hans Personlighed end ikke er Skygge
af Omskiftelse, hvordan saa end Lykken skifter i
Kunstens Verden; ja, den Begavede, der kan
gjøre dette, han er sandelig en Mester over
mange Mestere; han veed, hvad det vil sige,
% --- p. 70 ---
\setcounter{footnote}{0}%
\opage{70}at Gaven er mere end en Lykkens Gave; han
udtrykker personlig den Sandhed, at Gaven
kommer herovenfra, fra Lysenes Fader, hos
hvem der ikke er Forandring eller Skygge af
Omskiftelse.

Et Liv, der indvies til Kunsten, har altid
Mulighed til megen Lykke og megen Ulykke, til
bittre Smerter og til rige Nydelser; dette
gjælder ligefrem, forsaavidt Gaven opfattes som
Naturgave. Saaledes priisgivet for eventyrlig
Vexel af Held og Uheld, staaer Talentet ligefrem
under Skjæbnen.

Vil den Begavede derimod stille sig under
Forsynet, da er hans Liv, uden Hensyn til
Lykke og Ulykke, væsentlig indviet til ethisk Anstrengelse
og religiøs Lidelse. Og hvad er
saa alle Endelighedens Qvaklerier med Held
og Uheld, Anerkjendelse og Ikke-Anerkjendelse i
Sammenligning med den indre personlige Anstrengelse,
at skulle arbeide for et Maal, et
Ideal, der dog er et Middel, et Intet; at
skulle opflammes af Begeistring, og saa igjen
begeistres til at slukke Begeistringens Flamme!

Overveie vi nu dette, og betænke, at Kunstens
Børn ere Ideens Tjenere, udsendte til
% --- p. 71 ---
\setcounter{footnote}{0}%
\opage{71}Tjeneste for os; og betænke, hvad disse Tjenere
maae offre for vor Skyld, bære og lide for vor
Skyld, da have vi visselig ikke Aarsag til at
misunde dem deres aandige Fortrin; vi ere
meget mere forpligtede til at tjene dem igjen i
Alt, hvad vi kunne, komme dem imøde med Erkjendtlighed,
og ære dem som Forsynets Sendebud,
saasandt de virkelig (hvad vi jo maae troe)
i Alt bestræbe sig for at udtrykke Forsynets
Tanke.

See vi os om i Sandhedens Rige, da forefinde
vi to hinanden ganske modsatte Sphærer,
Videnskabens og Personlighedens. Den videnskabelige
Sandhed er i sit Væsen upersonlig;
den personlige Sandhed er i sit Væsen uvidenskabelig.

I Videnskaben gjælder Talentet; det er igjen
Endelighedens Forskjel imellem Begavede og
Ubegavede, der beviser, at Videnskaben tilhører
Verden. I Personlighedens Sphære er denne
Forskjel hævet. Sandheden siger: „Jeg priser
Dig, Fader, Himmelens og Jordens Herre!
at Du haver skjult dette for de Vise og Forstandige,
og aabenbaret det for de Umyndige“;
dette siger jo dog med andre Ord: „Jeg priser
% --- p. 72 ---
\setcounter{footnote}{0}%
\opage{72}Dig, Fader! at Du haver gjort den personlige
Sandhed uvidenskabelig.“

Begreberne Skjæbne og Forsyn ere selv
uvidenskabelige Begreber; deres Forskjel viser
sig først, idet Personligheden træder til, og indtager
sin Position.

Bøier Personligheden sig under den videnskabelige
Sandhed, saa seirer det Upersonlige,
og saa kommer Skjæbnen. Med denne Position
er det jo let nok at opløse den hele Troeslære
i Usandhed; men ved denne Position har Personligheden
tillige væsentlig opgivet sig selv, og
befinder sig i personlig Usandhed. Den videnskabelige
Sandhed har aldeles ingen Skyld i
denne Usandhed: Skylden ligger i det personlige
Selv.

Hævder Personligheden derimod sin oprindelige
Ret, og vender sig troende til Troens
Gjenstand: da er det en temmelig let Sag, at
blive de videnskabelige Indvendinger qvit; thi
de falde jo bort af sig selv.

Lad os da i Forsynets Navn anerkjende
vort religiøse Forhold til Gud som et personligt
Forhold; men lad os da ogsaa, i Forsynets
Navn, modstaae den Fristelse, at ville
% --- p. 73 ---
\setcounter{footnote}{0}%
\opage{73}have en „\emph{Troesvidenskab}“, eller med andre
Ord: at ville have dette Personlige igjen udlagt
i det Upersonlige.

Alt, hvad der videnskabeligt kan gjøres, for
at klare og sikkre Troens Opgaver, bestaaer i
en fuldstændig og tydelig Adskillelse af Videnskabens
og Personlighedens Enemærker. Ethvert Troesbegreb maa derfor fremstilles saaledes,
at det netop bliver kjendeligt ved dets
skarpe Modsætning til Vidensbegrebet, og Sphærernes
Ueensartethed paa alle Punkter efterviist
og belyst.

Saadanne Kategorier som: „religiøs Genius“,
„videnskabelig Naadegave“, høre til Indbildningens
Fostre, der hverken ere religiøse
eller videnskabelige.

Geniet er i dets umiddelbare Form en Naturens
Gave, og ingen Naturens Gave kan forholde
sig ligefrem understøttende til det Religiøse.

Det religiøse Øieblik er aldrig videnskabeligt;
det videnskabelige Øieblik er aldrig religiøst.

Det gaaer med Videnskabsmanden, som med
Kunstneren: vil han udtrykke Forsynets Tanke,
bliver han dragen i to modsatte Retninger.
Videnskaben er hans Maal, hans Værk, hans
% --- p. 74 ---
\setcounter{footnote}{0}%
\opage{74}verdslige Ideal. Arbeidende for dette Maal,
maa han bruge sin hele Bevidsthedskraft, for at
fjerne alt personlig Forstyrrende og glemme sig
selv i Interessen for den almene Sandhed.
Skal han saa tilbage i det Religiøse, da behøver
han igjen omvendt hele sin Bevidsthedskraft for
at samle sig personlig, vække sig selv til umiddelbar
Fornemmelse af Virkelighedens punktuelle
Indtryk, og erkjende, at det nu ikke gjælder
Dette eller Hiint i Almindelighed, men at det
gjælder ham, ham selv personlig for den personlige Gud.

Videnskabeligheden fremmer ikke det Religiøse,
den bedøver det Religiøse. Den, der
sidder i den dybeste Videnskabelighed, er ligesaa
langt fra det Religiøse, som Den, der ligger
i den dybeste Søvn.

Ogsaa heri viser det sig, at verdsligt Fortrin
er guddommelig Byrde. Det er langt lettere under Hverdagslivets Sysler, Sorger og
Glæder at arbeide inderlig i det Religiøse, end
at have det Religiøse i et Studerekammer. Du
kan bære paa en Møllesteen, og have det Religiøse;
men ingen Tænker paa Jorden kan have
det Religiøse i det Øieblik, han løser en viden%
% --- p. 75 ---
\setcounter{footnote}{0}%
\opage{75}skabelig Knude. En Kone kan spinde paa sin
Rok, og have det Religiøse; men en Mand,
var han end nok saa troende, kan ikke grunde
videnskabelig paa Guds Tilværelse, og til samme
Tid have det Religiøse.

At al Videnskab i og for sig, ligesom alt
Verdsligt i og for sig, fører bort fra Gud,
ligger deri, at al Videnskab fører bort fra det
Personlige. Naar det Personlige, det i Samvittighedsidealet
Punktuelle, først er behørig svækket,
kan Videnskabens Ideal da ogsaa let forvandle
sig til Guddommen selv. Derfor høres saadanne
Talemaader som: „Videnskaben fører os
Gud nærmere; Videnskabens Dyrkelse er i høiere
Forstand Guds Dyrkelse.“

Men er da Videnskaben ikke en god og fuldkommen
Gave? Jo, tilvisse! men det koster
utrolig Anstrengelse, naar den Videndes Existens
skal udtrykke den personlige Sandhed, at
Gaven kommer herovenfra, fra Lysenes Fader,
hos hvem der ikke er Forandring eller Skygge
af Omskiftelse.

Aandens Overlegenhed er Menneskets Lidelse:
Han, der haver den evige Aand, har lidt for
os Alle.

% --- p. 76 ---
\setcounter{footnote}{0}%
\opage{76}Overskue vi nu Fordelingen af de aandelige
Gaver i det Hele, da viser den sig, fra to forskjellige
Sider: verdslig uretfærdig, guddommelig
retfærdig.

Verdslighedens Aand, de relative Idealers,
de endelige Værkers og Bestræbelsers Aand,
er upersonlig, og just derfor uretfærdig. Naturens
Gaver uddeles i Flæng som Lykkens
Gaver: den Ene er begavet, den Anden indskrænket;
den Ene fuldfører et stort Værk, og
bliver noget Stort; den Anden udretter Intet,
og bliver i Ubetydelighed. Saa deles Arbeiderne:
den Begavede faaer det store Arbeide,
og er i Fordelen; den Ubegavede maa staae tilbage.
Saa strides de Arbeidende om Arbeidet;
saa opslides og fortæres de Arbeidende som
Midler, og Værkerne, de sidste Maal, blive
tilbage som Historiens Udbytte, som Babylons
Gruus, som Ægyptens Pyramider, som Mindesmærker
om de Masser af personlige Kræfter,
der i Tidernes Løb ere offrede til Skjæbnen.

Forsynets, Samvittighedsidealets, det evige
Formaals Aand er personlig, og just derfor
retfærdig. I dette Aandens Rige faaer hvert
Menneske sin Gave som personlig Opgave.
% --- p. 77 ---
\setcounter{footnote}{0}%
\opage{77}Den, der fik den største Kraft, fik ogsaa den
største Byrde. Her svarer Gaven nøiagtig til
Byrden, Værket til Ansvaret, Lønnen til Anstrengelsen;
her faaer den Største ikke Mere,
og kan ikke faae Mere, end den Mindste; thi
den Mindste saavel som den Største faaer jo,
enhver efter sin Anstrengelse i Troen, den hele
personlige Sandhed.

\stykke{IV.}{Det ene Fornødne.}

Men uden Tro er det umuligt at behage Gud; thi det
bør Den, som kommer frem for Gud, at troe, at
han er til, og at han bliver deres Belønner, som
søge ham.

„Naturen maa hjælpe sig selv“: denne Regel,
der vist har megen Gyldighed i Medicinen,
synes næsten at kunne overføres paa Ethiken,
naar man seer Alting verdslig. Thi er det
ikke forunderligt, hvorledes de personlige Naturanlæg
kunne spotte Opdragelse og Dannelse;
hvorledes en god Natur kan opdrage sig selv
trods slette Exempler, og en slet Natur forhærde
sig mod de bedste Exempler; hvorledes en selvstændig
Natur, snart bøielig, snart stridig, kan
% --- p. 78 ---
\setcounter{footnote}{0}%
\opage{78}opretholde fin Individualitet under stærke udvortes
Tryk, medens den svage, uselvstændige
Natur ikke ved nogensomhelst understøttende
Omgivelse kan hjælpes ud over sin medfødte
Umyndighed.

I verdslig Forstand er her naturligviis i
denne som i enhver anden timelig Ulighed en
Tilfældets Vilkaarlighed, en Skjæbnens Uretfærdighed;
men i guddommelig Forstand er
Skjæbnen Skyld, og Skylden er Menneskets
egen.

At dette forholder sig saaledes, kan jo
høres, naar de Paagjældende begynde at beklage
sig over Skjæbnen.

Klager den Svage, at han maa staae tilbage
for den Stærke; da kan han vel have Ret
i det Timelige; thi i alt Timeligt udretter den
stærke Natur langt mere end den svage. Men
for Gud anklager den Svage sig selv; thi for
Gud gjælder Troen som det ene Fornødne, og
i Troen er Forskjellen hævet.

For at troe, maa den Stærke blive svag,
som den Svageste; thi Gud er kun stærk i de
Svage. Skulde her være nogen endelig Forskjel,
da maatte Fordelen jo snarere være paa
% --- p. 79 ---
\setcounter{footnote}{0}%
\opage{79}den Svageres Side; men her er ingen Forskjel:
den naturlige Svaghed og den naturlige
Styrke forsvinder ligelig i den Svaghed, som
ene gjælder for Gud. Saa sættes Ligheden for
Gud i det Indre, og al den Ulighed, der tilhører
det Ydre, den hele udadgaaende Virksomhed,
er da at betragte som det Forsvindende,
der ikke vedrører det personlige Selv.

Er det et uroligt, lidenskabeligt, i indre
Sønderrevethed martret Sind, der vil sammenligne
sig med den besindige, rolige, af Naturen
harmonisk anlagte Individualitet, og beklage sig
over sin Skjæbne: da vil ogsaa denne Klage
høres af Gud som en Selvanklage; thi Troen
fordres, og i Troen er Forskjellen hævet.

Den af Naturen Rolige maa rystes og uroliges
til det Yderste, for at naae Troen; og
den af Naturen Urolige finder jo Ro og Fred,
naar han kommer til Troen. Troen er ikke
lidt Ro og lidt Uro, eller lutter Fred i Nogle
og lutter Ufred i Andre; men Troen er overalt
og i Alle det Eviges Uro i al naturlig Ro,
det Eviges Fred i al naturlig Ufred. Den
Ulighed i Blandingen af Fred og Ufred; der
som en Følge af Individualiteternes Naturfor%
% --- p. 80 ---
\setcounter{footnote}{0}%
\opage{80}skjellighed bliver tilbage, tilhører den timelige
Verden, den timelige Kamp, og er det personlige
Selv underordnet.

Vil den blidere, blødere Natur beklage sig
i Sammenligning med den haardere og haardføre, at denne har langt lettere ved at bære
Livets Tilskikkelser; da høres ogsaa denne Klagende
af Gud som En, der anklager sig selv.
Tilværelsen kan være haard og streng; den af
Naturen Haardføre har i det Timelige det
umiskjendelige Fortrin, at han lettere hærder
sig mod Skjæbnens Slag, end den Ulykkelige,
der fik det blødere, blidere Sind. Men i Troen
er Forskjellen hævet. Troens Opgave er Ingen
for streng, thi den er uendelig mild; den
er Ingen for mild, thi den er uendelig streng.
Den himmelske Fader er ikke som en jordisk
Fader lidt mere føielig mod de Føielige, lidt
mere haard imod de Haarde; nei, Forsynets
Gud er væsentlig lige naadig og lige streng
imod alle Naturer.

Som der kan være en Gjenstridighed i
Haardhed, saa kan der ogsaa være en Gjenstridighed i klynkende Blødagtighed: den hellige
Strenghed fortærer al naturlig Gjenstridighed;
% --- p. 81 ---
\setcounter{footnote}{0}%
\opage{81}men, naar den guddommelige Ild har fortæret
det menneskelig Gjenstridige, forvandler den sig
til Mildhed, til uendelig Barmhjertighed. Derfor
skal intet Menneske kjende sin Tro derpaa, at
han har Strengheden alene, og Ingen derpaa,
at han har Mildheden alene; men under alle
Individualiteternes Forskjelligheder kan Enhver
kjende sin Tro derpaa, at den altid har Strengheden
i Mildhed og Mildheden i Strenghed.

Stundom klage vi Mennesker over, at Troen
er for overjordisk, for aandelig, for overspændt,
saa at den ikke ret tillader os at see paa det
Timelige, glæde os over det reent Menneskelige
og virke i det Nærværende. Det er naturligt,
at denne Klage opkommer hos Enhver,
der i nogen Grad anstrenger sig med Troens
Opgave; lad os da klage, den Ene over Hiint,
den Anden over Dette, naar vi kun vide, at
vi for Gud anklage os selv.

Eet er fornødent; dette Ene er Troen, i
Troen ubetinget Anerkjendelse af Samvittighedens
Fordringer, i Troen ubetinget Ydmygelse
under Personlighedens Ideal, i Troen
ubetinget Lydighed mod Idealets Love, i Troen
ubetinget Tillid til Idealets Forjættelser; naar
% --- p. 82 ---
\setcounter{footnote}{0}%
\opage{82}et Menneske i inderlig Alvor stræber at fastholde
dette Ene: da leve han i Verden, som
han kan og vil, eensom eller i Selskab, fastende
i Ørkenen eller nydende i Samfundet: der skal
være Velsignelse i det Altsammen, og det skal
Altsammen komme ham tilgode.

Men uden Tro er det umuligt.

Det Timelige som timeligt er ikke til en vis
Grad ligegyldigt for Gud; det er uendelig ligegyldigt.
Timeligt Savn, timelig Smerte, timelig
Glæde, timelig Nydelse, det er i sig Altsammen
Intet; her gjælder ingen Endelighedens
Medynk, ingen Klynken og ingen Jamren, alt
Timeligt gaaer og skal gaae efter Naturens
Orden: Mennesket i Kiød og Blod er „Støv
og Aske“, et Intet for Gud, som Verden i al
dens Herlighed og al dens Elendighed er et
Intet for Gud, den Almægtige. Men, just fordi
det Timelige i dets Timelighed er uendelig ligegyldigt,
er det som Middel og Betingelse for
Menneskets personlige Opdragelse af uendelig
Gyldighed. Forsynet har ikke som Odin kun
eet Øie; Forsynet har Millioner Øine, et
Øie med hvert Menneske, et Øie i hver Samvittighed.
I Slægternes Mylder, i Begiven%
% --- p. 83 ---
\setcounter{footnote}{0}%
\opage{83}hedernes Trængsel er ingen Sjæl forglemt, og
intet Tilfælde tilfældigt.

Verdens Magter larme, den Almægtige tier;
Verdens Magter ruste sig og forsvare deres
Ret, den Almægtige er saa stille og rolig!
Lad dem, der kunne nøies med lidt Ret og
lidt Uret, lidt Skyld og lidt Skjæbne, udtyde
dette, som de ville: hvert Menneske, der i Samvittighedens
Navn vil frelses af Skjæbnens
Vold, maa erkjende, at Forsynet er aabenbaret,
at Knuden er løst, at der er Hjælp for Enhver,
der vil lade sig hjælpe.

Men uden Tro er det umuligt; thi Troen
paa Forsynet er Forsynets Hjælp.

\begin{center}---\end{center}

\end{document}
